% generated by GAPDoc2LaTeX from XML source (Frank Luebeck)
\documentclass[a4paper,11pt]{report}

\usepackage[top=37mm,bottom=37mm,left=27mm,right=27mm]{geometry}
\sloppy
\pagestyle{myheadings}
\usepackage{amssymb}
\usepackage[utf8]{inputenc}
\usepackage{makeidx}
\makeindex
\usepackage{color}
\definecolor{FireBrick}{rgb}{0.5812,0.0074,0.0083}
\definecolor{RoyalBlue}{rgb}{0.0236,0.0894,0.6179}
\definecolor{RoyalGreen}{rgb}{0.0236,0.6179,0.0894}
\definecolor{RoyalRed}{rgb}{0.6179,0.0236,0.0894}
\definecolor{LightBlue}{rgb}{0.8544,0.9511,1.0000}
\definecolor{Black}{rgb}{0.0,0.0,0.0}

\definecolor{linkColor}{rgb}{0.0,0.0,0.554}
\definecolor{citeColor}{rgb}{0.0,0.0,0.554}
\definecolor{fileColor}{rgb}{0.0,0.0,0.554}
\definecolor{urlColor}{rgb}{0.0,0.0,0.554}
\definecolor{promptColor}{rgb}{0.0,0.0,0.589}
\definecolor{brkpromptColor}{rgb}{0.589,0.0,0.0}
\definecolor{gapinputColor}{rgb}{0.589,0.0,0.0}
\definecolor{gapoutputColor}{rgb}{0.0,0.0,0.0}

%%  for a long time these were red and blue by default,
%%  now black, but keep variables to overwrite
\definecolor{FuncColor}{rgb}{0.0,0.0,0.0}
%% strange name because of pdflatex bug:
\definecolor{Chapter }{rgb}{0.0,0.0,0.0}
\definecolor{DarkOlive}{rgb}{0.1047,0.2412,0.0064}


\usepackage{fancyvrb}

\usepackage{mathptmx,helvet}
\usepackage[T1]{fontenc}
\usepackage{textcomp}


\usepackage[
            pdftex=true,
            bookmarks=true,        
            a4paper=true,
            pdftitle={Written with GAPDoc},
            pdfcreator={LaTeX with hyperref package / GAPDoc},
            colorlinks=true,
            backref=page,
            breaklinks=true,
            linkcolor=linkColor,
            citecolor=citeColor,
            filecolor=fileColor,
            urlcolor=urlColor,
            pdfpagemode={UseNone}, 
           ]{hyperref}

\newcommand{\maintitlesize}{\fontsize{50}{55}\selectfont}

% write page numbers to a .pnr log file for online help
\newwrite\pagenrlog
\immediate\openout\pagenrlog =\jobname.pnr
\immediate\write\pagenrlog{PAGENRS := [}
\newcommand{\logpage}[1]{\protect\write\pagenrlog{#1, \thepage,}}
%% were never documented, give conflicts with some additional packages

\newcommand{\GAP}{\textsf{GAP}}

%% nicer description environments, allows long labels
\usepackage{enumitem}
\setdescription{style=nextline}

%% depth of toc
\setcounter{tocdepth}{1}





%% command for ColorPrompt style examples
\newcommand{\gapprompt}[1]{\color{promptColor}{\bfseries #1}}
\newcommand{\gapbrkprompt}[1]{\color{brkpromptColor}{\bfseries #1}}
\newcommand{\gapinput}[1]{\color{gapinputColor}{#1}}


\begin{document}

\logpage{[ 0, 0, 0 ]}
\begin{titlepage}
\mbox{}\vfill

\begin{center}{\maintitlesize \textbf{ NormalizInterface \mbox{}}}\\
\vfill

\hypersetup{pdftitle= NormalizInterface }
\markright{\scriptsize \mbox{}\hfill  NormalizInterface  \hfill\mbox{}}
{\Huge \textbf{ \textsf{GAP} wrapper for Normaliz \mbox{}}}\\
\vfill

{\Huge  1.5.0 \mbox{}}\\[1cm]
{ 24 March 2026 \mbox{}}\\[1cm]
\mbox{}\\[2cm]
{\Large \textbf{ Sebastian Gutsche\\
  \mbox{}}}\\
{\Large \textbf{ Max Horn\\
    \mbox{}}}\\
{\Large \textbf{ Christof S{\"o}ger\\
  \mbox{}}}\\
\hypersetup{pdfauthor= Sebastian Gutsche\\
  ;  Max Horn\\
    ;  Christof S{\"o}ger\\
  }
\end{center}\vfill

\mbox{}\\
{\mbox{}\\
\small \noindent \textbf{ Sebastian Gutsche\\
  }  Email: \href{mailto://gutsche@mathematik.uni-siegen.de} {\texttt{gutsche@mathematik.uni\texttt{\symbol{45}}siegen.de}}}\\
{\mbox{}\\
\small \noindent \textbf{ Max Horn\\
    }  Email: \href{mailto://mhorn@rptu.de} {\texttt{mhorn@rptu.de}}\\
  Homepage: \href{https://www.quendi.de/math} {\texttt{https://www.quendi.de/math}}\\
  Address: \begin{minipage}[t]{8cm}\noindent
 Fachbereich Mathematik\\
 RPTU Kaiserslautern\texttt{\symbol{45}}Landau\\
 Gottlieb\texttt{\symbol{45}}Daimler\texttt{\symbol{45}}Stra{\ss}e 48\\
 67663 Kaiserslautern\\
 Germany\\
 \end{minipage}
}\\
{\mbox{}\\
\small \noindent \textbf{ Christof S{\"o}ger\\
  }  Email: \href{mailto://csoeger@uos.de} {\texttt{csoeger@uos.de}}}\\
\end{titlepage}

\newpage\setcounter{page}{2}
\newpage

\def\contentsname{Contents\logpage{[ 0, 0, 1 ]}}

\tableofcontents
\newpage

     
\chapter{\textcolor{Chapter }{Introduction}}\label{Chapter_Introduction}
\logpage{[ 1, 0, 0 ]}
\hyperdef{L}{X7DFB63A97E67C0A1}{}
{
  

 
\section{\textcolor{Chapter }{What is the purpose of the this package?}}\label{Chapter_Introduction_Section_What_is_the_purpose_of_the_this_package}
\logpage{[ 1, 1, 0 ]}
\hyperdef{L}{X86B6D8F583923CB5}{}
{
  

 Normaliz \cite{normaliz} is a software for computations with rational cones and affine monoids. It
pursues two main computational goals: finding the Hilbert basis, a minimal
generating system of the monoid of lattice points of a cone; and counting
elements degree\texttt{\symbol{45}}wise in a generating function, the Hilbert
series. As a recent extension, Normaliz can handle unbounded polyhedra. The
Hilbert basis computation can be considered as solving a linear diophantine
system of inhomogeneous equations, inequalities and congruences. 

 

 This package encapsulates a libnormaliz cone and gives access to it in the \textsf{GAP} environment. In this way \textsf{GAP} can be used as interactive interface to libnormaliz. 

 

 }

 }

   
\chapter{\textcolor{Chapter }{Functions}}\label{Chapter_Functions}
\logpage{[ 2, 0, 0 ]}
\hyperdef{L}{X86FA580F8055B274}{}
{
  

 In this chapter we describe the functions offered by \mbox{\texttt{\mdseries\slshape NormalizInterface}}. All functions supplied by this package start with ``Nmz''. For examples see the chapter \hyperref[Chapter_Examples]{`Examples'}. 

 
\section{\textcolor{Chapter }{Create a NmzCone}}\label{Chapter_Functions_Section_Create_a_NmzCone}
\logpage{[ 2, 1, 0 ]}
\hyperdef{L}{X788BCD4B87853B34}{}
{
  

 To create a cone object use \texttt{NmzCone} (\ref{NmzCone}). 

 

\subsection{\textcolor{Chapter }{NmzCone}}
\logpage{[ 2, 1, 1 ]}\nobreak
\hyperdef{L}{X7DBDE20A8515FF5D}{}
{\noindent\textcolor{FuncColor}{$\triangleright$\enspace\texttt{NmzCone({\mdseries\slshape list})\index{NmzCone@\texttt{NmzCone}}
\label{NmzCone}
}\hfill{\scriptsize (function)}}\\
\textbf{\indent Returns:\ }
NmzCone 



 Creates a NmzCone. The \mbox{\texttt{\mdseries\slshape list}} argument should contain an even number of elements, alternating between a
string and a integer matrix. The string has to correspond to a Normaliz input
type string and the following matrix will be interpreted as input of that
type. 

 See the Normaliz manual for the Normaliz version loaded by your version of
NormalizInterface for a detailed description of which input type strings are
supported and what arguments they take. 

 
\begin{Verbatim}[commandchars=!@|,fontsize=\small,frame=single,label=Example]
  !gapprompt@gap>| !gapinput@cone := NmzCone(["integral_closure",[[2,1],[1,3]]]);|
  <a Normaliz cone>
\end{Verbatim}
 }

 }

 
\section{\textcolor{Chapter }{Use a NmzCone}}\label{Chapter_Functions_Section_Use_a_NmzCone}
\logpage{[ 2, 2, 0 ]}
\hyperdef{L}{X7D0D9BF983F5DABB}{}
{
  

 

\subsection{\textcolor{Chapter }{NmzHasConeProperty}}
\logpage{[ 2, 2, 1 ]}\nobreak
\hyperdef{L}{X87E0E6967E47394E}{}
{\noindent\textcolor{FuncColor}{$\triangleright$\enspace\texttt{NmzHasConeProperty({\mdseries\slshape cone, property})\index{NmzHasConeProperty@\texttt{NmzHasConeProperty}}
\label{NmzHasConeProperty}
}\hfill{\scriptsize (function)}}\\
\textbf{\indent Returns:\ }
whether the cone has already computed the given property 



 See \texttt{NmzConeProperty} (\ref{NmzConeProperty}) for a list of recognized properties. 

 
\begin{Verbatim}[commandchars=!@|,fontsize=\small,frame=single,label=Example]
  !gapprompt@gap>| !gapinput@NmzHasConeProperty(cone, "ExtremeRays");|
  false
\end{Verbatim}
 }

 

\subsection{\textcolor{Chapter }{NmzKnownConeProperties}}
\logpage{[ 2, 2, 2 ]}\nobreak
\hyperdef{L}{X8488EAA478AA3706}{}
{\noindent\textcolor{FuncColor}{$\triangleright$\enspace\texttt{NmzKnownConeProperties({\mdseries\slshape cone})\index{NmzKnownConeProperties@\texttt{NmzKnownConeProperties}}
\label{NmzKnownConeProperties}
}\hfill{\scriptsize (function)}}\\
\textbf{\indent Returns:\ }
a list of strings representing the known (computed) cone properties 



 Given a Normaliz cone object, return a list of all properties already computed
for the cone. 

 
\begin{Verbatim}[commandchars=!@|,fontsize=\small,frame=single,label=Example]
  !gapprompt@gap>| !gapinput@NmzKnownConeProperties(cone);|
  [  ]
\end{Verbatim}
 }

 

\subsection{\textcolor{Chapter }{NmzSetVerboseDefault}}
\logpage{[ 2, 2, 3 ]}\nobreak
\hyperdef{L}{X7AD4C81887308070}{}
{\noindent\textcolor{FuncColor}{$\triangleright$\enspace\texttt{NmzSetVerboseDefault({\mdseries\slshape verboseFlag})\index{NmzSetVerboseDefault@\texttt{NmzSetVerboseDefault}}
\label{NmzSetVerboseDefault}
}\hfill{\scriptsize (function)}}\\
\textbf{\indent Returns:\ }
the previous verbosity 



 Set the global default verbosity state in libnormaliz. This will influence all
NmzCone created afterwards, but not any existing ones. 

 See also \texttt{NmzSetVerbose} (\ref{NmzSetVerbose}) }

 

\subsection{\textcolor{Chapter }{NmzSetVerbose}}
\logpage{[ 2, 2, 4 ]}\nobreak
\hyperdef{L}{X84A5129078C5D198}{}
{\noindent\textcolor{FuncColor}{$\triangleright$\enspace\texttt{NmzSetVerbose({\mdseries\slshape cone, verboseFlag})\index{NmzSetVerbose@\texttt{NmzSetVerbose}}
\label{NmzSetVerbose}
}\hfill{\scriptsize (function)}}\\
\textbf{\indent Returns:\ }
the previous verbosity 



 Set the verbosity state for a cone. 

 See also \texttt{NmzSetVerboseDefault} (\ref{NmzSetVerboseDefault}) }

 

\subsection{\textcolor{Chapter }{NmzCompute}}
\logpage{[ 2, 2, 5 ]}\nobreak
\hyperdef{L}{X82752899822ED2F0}{}
{\noindent\textcolor{FuncColor}{$\triangleright$\enspace\texttt{NmzCompute({\mdseries\slshape cone[, propnames]})\index{NmzCompute@\texttt{NmzCompute}}
\label{NmzCompute}
}\hfill{\scriptsize (function)}}\\
\textbf{\indent Returns:\ }
a boolean indicating success 



 Start computing properties of the given cone. The first parameter indicates a
cone object, the second parameter is either a single string, or a list of
strings, which indicate what should be computed. 

 

 The single parameter version is equivalent to \texttt{NmzCone(cone, ["DefaultMode"])}. See \texttt{NmzConeProperty} (\ref{NmzConeProperty}) for a list of recognized properties. 
\begin{Verbatim}[commandchars=!@|,fontsize=\small,frame=single,label=Example]
  !gapprompt@gap>| !gapinput@NmzKnownConeProperties(cone);|
  [  ]
  !gapprompt@gap>| !gapinput@NmzCompute(cone, ["SupportHyperplanes", "IsPointed"]);|
  true
  !gapprompt@gap>| !gapinput@NmzKnownConeProperties(cone);|
  [ "EmbeddingDim", "ExtremeRays", "Generators", "InternalIndex", 
    "IsDeg1ExtremeRays", "IsInhomogeneous", "IsPointed", "MaximalSubspace", 
    "OriginalMonoidGenerators", "Rank", "Sublattice", "SupportHyperplanes" ]
  !gapprompt@gap>| !gapinput@NmzCompute(cone);;|
  !gapprompt@gap>| !gapinput@NmzKnownConeProperties(cone);|
  [ "ClassGroup", "EmbeddingDim", "ExtremeRays", "Generators", "HilbertBasis", 
    "InternalIndex", "IsDeg1ExtremeRays", "IsInhomogeneous", 
    "IsIntegrallyClosed", "IsPointed", "IsTriangulationNested", 
    "IsTriangulationPartial", "MaximalSubspace", "OriginalMonoidGenerators", 
    "Rank", "Sublattice", "SupportHyperplanes", "TriangulationDetSum", 
    "TriangulationSize", "UnitGroupIndex" ]
\end{Verbatim}
 }

 

\subsection{\textcolor{Chapter }{NmzConeProperty}}
\logpage{[ 2, 2, 6 ]}\nobreak
\hyperdef{L}{X78ED435078E3E377}{}
{\noindent\textcolor{FuncColor}{$\triangleright$\enspace\texttt{NmzConeProperty({\mdseries\slshape cone, property})\index{NmzConeProperty@\texttt{NmzConeProperty}}
\label{NmzConeProperty}
}\hfill{\scriptsize (function)}}\\
\textbf{\indent Returns:\ }
the result of the computation, type depends on the property 



 Triggers the computation of the property of the cone and returns the result.
If the property was already known, it is not recomputed. Currently the
following strings are recognized as properties: 
\begin{itemize}
\item \texttt{Generators} see \texttt{NmzGenerators} (\ref{NmzGenerators}),
\item \texttt{ExtremeRays} see \texttt{NmzExtremeRays} (\ref{NmzExtremeRays}),
\item \texttt{VerticesOfPolyhedron} see \texttt{NmzVerticesOfPolyhedron} (\ref{NmzVerticesOfPolyhedron}),
\item \texttt{SupportHyperplanes} see \texttt{NmzSupportHyperplanes} (\ref{NmzSupportHyperplanes}),
\item \texttt{TriangulationSize} see \texttt{NmzTriangulationSize} (\ref{NmzTriangulationSize}),
\item \texttt{TriangulationDetSum} see \texttt{NmzTriangulationDetSum} (\ref{NmzTriangulationDetSum}),
\item \texttt{Triangulation} see \texttt{NmzTriangulation} (\ref{NmzTriangulation}),
\item \texttt{Multiplicity} see \texttt{NmzMultiplicity} (\ref{NmzMultiplicity}),
\item \texttt{RecessionRank} see \texttt{NmzRecessionRank} (\ref{NmzRecessionRank}),
\item \texttt{AffineDim} see \texttt{NmzAffineDim} (\ref{NmzAffineDim}),
\item \texttt{ModuleRank} see \texttt{NmzModuleRank} (\ref{NmzModuleRank}),
\item \texttt{HilbertBasis} see \texttt{NmzHilbertBasis} (\ref{NmzHilbertBasis}),
\item \texttt{ModuleGenerators} see \texttt{NmzModuleGenerators} (\ref{NmzModuleGenerators}),
\item \texttt{Deg1Elements} see \texttt{NmzDeg1Elements} (\ref{NmzDeg1Elements}),
\item \texttt{HilbertSeries} see \texttt{NmzHilbertSeries} (\ref{NmzHilbertSeries}),
\item \texttt{HilbertQuasiPolynomial} see \texttt{NmzHilbertQuasiPolynomial} (\ref{NmzHilbertQuasiPolynomial}),
\item \texttt{Grading} see \texttt{NmzGrading} (\ref{NmzGrading}),
\item \texttt{IsPointed} see \texttt{NmzIsPointed} (\ref{NmzIsPointed}),
\item \texttt{IsDeg1ExtremeRays} see \texttt{NmzIsDeg1ExtremeRays} (\ref{NmzIsDeg1ExtremeRays}),
\item \texttt{IsDeg1HilbertBasis} see \texttt{NmzIsDeg1HilbertBasis} (\ref{NmzIsDeg1HilbertBasis}),
\item \texttt{IsIntegrallyClosed} see \texttt{NmzIsIntegrallyClosed} (\ref{NmzIsIntegrallyClosed}),
\item \texttt{OriginalMonoidGenerators} see \texttt{NmzOriginalMonoidGenerators} (\ref{NmzOriginalMonoidGenerators}),
\item \texttt{IsReesPrimary} see \texttt{NmzIsReesPrimary} (\ref{NmzIsReesPrimary}),
\item \texttt{ReesPrimaryMultiplicity} see \texttt{NmzReesPrimaryMultiplicity} (\ref{NmzReesPrimaryMultiplicity}),
\item \texttt{ExcludedFaces} see \texttt{NmzExcludedFaces} (\ref{NmzExcludedFaces}),
\item \texttt{Dehomogenization} see \texttt{NmzDehomogenization} (\ref{NmzDehomogenization}),
\item \texttt{InclusionExclusionData} see \texttt{NmzInclusionExclusionData} (\ref{NmzInclusionExclusionData}),
\item \texttt{ClassGroup} see \texttt{NmzClassGroup} (\ref{NmzClassGroup}),
\item \texttt{ModuleGeneratorsOverOriginalMonoid} see \texttt{NmzModuleGeneratorsOverOriginalMonoid} (\ref{NmzModuleGeneratorsOverOriginalMonoid}),
\item \texttt{Sublattice} computes the efficient sublattice and returns a bool signaling whether the
computation was successful. Actual data connected to it can be accessed by \texttt{NmzRank} (\ref{NmzRank}), \texttt{NmzEquations} (\ref{NmzEquations}), \texttt{NmzCongruences} (\ref{NmzCongruences}), and \texttt{NmzBasisChange} (\ref{NmzBasisChange}).
\end{itemize}
 

 Additionally also the following compute options are accepted as property. They
modify what and how should be computed, and return True after a successful
computation. 
\begin{itemize}
\item \texttt{Approximate} approximate the rational polytope by an integral polytope, currently only
useful in combination with \texttt{Deg1Elements}.
\item \texttt{BottomDecomposition} use the best possible triangulation (with respect to the sum of determinants)
using the given generators.
\item \texttt{DefaultMode} try to compute what is possible and do not throw an exception when something
cannot be computed.
\item \texttt{DualMode}activates the dual algorithm for the computation of the Hilbert basis and
degree 1 elements. Includes \texttt{HilbertBasis}, unless \texttt{Deg1Elements} is set. Often a good choice if you start from constraints.
\item \texttt{KeepOrder} forbids to reorder the generators. Blocks \texttt{BottomDecomposition}.
\end{itemize}
 

 All the properties above can be given to \texttt{NmzCompute} (\ref{NmzCompute}). There you can combine different properties, e.g. give some properties that
you would like to know and add some compute options. 

 See the Normaliz manual for a detailed description. 

 }

 

\subsection{\textcolor{Chapter }{NmzPrintConeProperties}}
\logpage{[ 2, 2, 7 ]}\nobreak
\hyperdef{L}{X87C644A6861CFCD9}{}
{\noindent\textcolor{FuncColor}{$\triangleright$\enspace\texttt{NmzPrintConeProperties({\mdseries\slshape cone})\index{NmzPrintConeProperties@\texttt{NmzPrintConeProperties}}
\label{NmzPrintConeProperties}
}\hfill{\scriptsize (function)}}\\


 Print an overview of all known properties of the given cone, as well as their
values. }

 }

 
\section{\textcolor{Chapter }{Cone properties}}\label{Chapter_Functions_Section_Cone_properties}
\logpage{[ 2, 3, 0 ]}
\hyperdef{L}{X81EC46D08211949E}{}
{
  

 

\subsection{\textcolor{Chapter }{NmzAffineDim}}
\logpage{[ 2, 3, 1 ]}\nobreak
\hyperdef{L}{X79232A2D7B1CAFC2}{}
{\noindent\textcolor{FuncColor}{$\triangleright$\enspace\texttt{NmzAffineDim({\mdseries\slshape cone})\index{NmzAffineDim@\texttt{NmzAffineDim}}
\label{NmzAffineDim}
}\hfill{\scriptsize (function)}}\\
\textbf{\indent Returns:\ }
the affine dimension 



 

 The affine dimension of the polyhedron in inhomogeneous computations. Its
computation is triggered if necessary. 

 This is an alias for \texttt{NmzConeProperty( cone, "AffineDim" );} see \texttt{NmzConeProperty} (\ref{NmzConeProperty}). }

 

\subsection{\textcolor{Chapter }{NmzAllGeneratorsTriangulation}}
\logpage{[ 2, 3, 2 ]}\nobreak
\hyperdef{L}{X872772CB78B4E2E8}{}
{\noindent\textcolor{FuncColor}{$\triangleright$\enspace\texttt{NmzAllGeneratorsTriangulation({\mdseries\slshape cone})\index{NmzAllGeneratorsTriangulation@\texttt{NmzAllGeneratorsTriangulation}}
\label{NmzAllGeneratorsTriangulation}
}\hfill{\scriptsize (function)}}\\


 This is an alias for \texttt{NmzConeProperty( cone, "AllGeneratorsTriangulation" );} see \texttt{NmzConeProperty} (\ref{NmzConeProperty}). }

 

\subsection{\textcolor{Chapter }{NmzAmbientAutomorphisms}}
\logpage{[ 2, 3, 3 ]}\nobreak
\hyperdef{L}{X7CCFE4FE84145C13}{}
{\noindent\textcolor{FuncColor}{$\triangleright$\enspace\texttt{NmzAmbientAutomorphisms({\mdseries\slshape cone})\index{NmzAmbientAutomorphisms@\texttt{NmzAmbientAutomorphisms}}
\label{NmzAmbientAutomorphisms}
}\hfill{\scriptsize (function)}}\\


 This is an alias for \texttt{NmzConeProperty( cone, "AmbientAutomorphisms" );} see \texttt{NmzConeProperty} (\ref{NmzConeProperty}). }

 

\subsection{\textcolor{Chapter }{NmzApproximate}}
\logpage{[ 2, 3, 4 ]}\nobreak
\hyperdef{L}{X7D0B2A4978D88CD9}{}
{\noindent\textcolor{FuncColor}{$\triangleright$\enspace\texttt{NmzApproximate({\mdseries\slshape cone})\index{NmzApproximate@\texttt{NmzApproximate}}
\label{NmzApproximate}
}\hfill{\scriptsize (function)}}\\


 This is an alias for \texttt{NmzConeProperty( cone, "Approximate" );} see \texttt{NmzConeProperty} (\ref{NmzConeProperty}). }

 

\subsection{\textcolor{Chapter }{NmzAutomorphisms}}
\logpage{[ 2, 3, 5 ]}\nobreak
\hyperdef{L}{X78CB11547A7EF6C0}{}
{\noindent\textcolor{FuncColor}{$\triangleright$\enspace\texttt{NmzAutomorphisms({\mdseries\slshape cone})\index{NmzAutomorphisms@\texttt{NmzAutomorphisms}}
\label{NmzAutomorphisms}
}\hfill{\scriptsize (function)}}\\


 This is an alias for \texttt{NmzConeProperty( cone, "Automorphisms" );} see \texttt{NmzConeProperty} (\ref{NmzConeProperty}). }

 

\subsection{\textcolor{Chapter }{NmzAxesScaling}}
\logpage{[ 2, 3, 6 ]}\nobreak
\hyperdef{L}{X7F80B01C7F6760B9}{}
{\noindent\textcolor{FuncColor}{$\triangleright$\enspace\texttt{NmzAxesScaling({\mdseries\slshape cone})\index{NmzAxesScaling@\texttt{NmzAxesScaling}}
\label{NmzAxesScaling}
}\hfill{\scriptsize (function)}}\\


 This is an alias for \texttt{NmzConeProperty( cone, "AxesScaling" );} see \texttt{NmzConeProperty} (\ref{NmzConeProperty}). }

 

\subsection{\textcolor{Chapter }{NmzBasicStanleyDec}}
\logpage{[ 2, 3, 7 ]}\nobreak
\hyperdef{L}{X802E37267CCB35EB}{}
{\noindent\textcolor{FuncColor}{$\triangleright$\enspace\texttt{NmzBasicStanleyDec({\mdseries\slshape cone})\index{NmzBasicStanleyDec@\texttt{NmzBasicStanleyDec}}
\label{NmzBasicStanleyDec}
}\hfill{\scriptsize (function)}}\\


 This is an alias for \texttt{NmzConeProperty( cone, "BasicStanleyDec" );} see \texttt{NmzConeProperty} (\ref{NmzConeProperty}). }

 

\subsection{\textcolor{Chapter }{NmzBasicTriangulation}}
\logpage{[ 2, 3, 8 ]}\nobreak
\hyperdef{L}{X7C49ED54861542CF}{}
{\noindent\textcolor{FuncColor}{$\triangleright$\enspace\texttt{NmzBasicTriangulation({\mdseries\slshape cone})\index{NmzBasicTriangulation@\texttt{NmzBasicTriangulation}}
\label{NmzBasicTriangulation}
}\hfill{\scriptsize (function)}}\\


 This is an alias for \texttt{NmzConeProperty( cone, "BasicTriangulation" );} see \texttt{NmzConeProperty} (\ref{NmzConeProperty}). }

 

\subsection{\textcolor{Chapter }{NmzBigInt}}
\logpage{[ 2, 3, 9 ]}\nobreak
\hyperdef{L}{X7C1A8AE178FC3315}{}
{\noindent\textcolor{FuncColor}{$\triangleright$\enspace\texttt{NmzBigInt({\mdseries\slshape cone})\index{NmzBigInt@\texttt{NmzBigInt}}
\label{NmzBigInt}
}\hfill{\scriptsize (function)}}\\


 This is an alias for \texttt{NmzConeProperty( cone, "BigInt" );} see \texttt{NmzConeProperty} (\ref{NmzConeProperty}). }

 

\subsection{\textcolor{Chapter }{NmzBottomDecomposition}}
\logpage{[ 2, 3, 10 ]}\nobreak
\hyperdef{L}{X83485BC4780CA904}{}
{\noindent\textcolor{FuncColor}{$\triangleright$\enspace\texttt{NmzBottomDecomposition({\mdseries\slshape cone})\index{NmzBottomDecomposition@\texttt{NmzBottomDecomposition}}
\label{NmzBottomDecomposition}
}\hfill{\scriptsize (function)}}\\


 This is an alias for \texttt{NmzConeProperty( cone, "BottomDecomposition" );} see \texttt{NmzConeProperty} (\ref{NmzConeProperty}). }

 

\subsection{\textcolor{Chapter }{NmzClassGroup}}
\logpage{[ 2, 3, 11 ]}\nobreak
\hyperdef{L}{X7B12EE3D7ED4F9EB}{}
{\noindent\textcolor{FuncColor}{$\triangleright$\enspace\texttt{NmzClassGroup({\mdseries\slshape cone})\index{NmzClassGroup@\texttt{NmzClassGroup}}
\label{NmzClassGroup}
}\hfill{\scriptsize (function)}}\\
\textbf{\indent Returns:\ }
the class group in a special format 



 

 A normal affine monoid $M$ has a well\texttt{\symbol{45}}defined divisor class group. It is naturally
isomorphic to the divisor class group of $K[M]$ where $K$ is a field (or any unique factorization domain). We represent it as a vector
where the first entry is the rank. It is followed by sequence of pairs of
entries $n,m$. Such two entries represent a free cyclic summand $(\mathbb{Z}/n\mathbb{Z})^m$. Not allowed in inhomogeneous computations. 

 This is an alias for \texttt{NmzConeProperty( cone, "ClassGroup" );} see \texttt{NmzConeProperty} (\ref{NmzConeProperty}). }

 

\subsection{\textcolor{Chapter }{NmzCombinatorialAutomorphisms}}
\logpage{[ 2, 3, 12 ]}\nobreak
\hyperdef{L}{X7E49A40886FF88D1}{}
{\noindent\textcolor{FuncColor}{$\triangleright$\enspace\texttt{NmzCombinatorialAutomorphisms({\mdseries\slshape cone})\index{NmzCombinatorialAutomorphisms@\texttt{NmzCombinatorialAutomorphisms}}
\label{NmzCombinatorialAutomorphisms}
}\hfill{\scriptsize (function)}}\\


 This is an alias for \texttt{NmzConeProperty( cone, "CombinatorialAutomorphisms" );} see \texttt{NmzConeProperty} (\ref{NmzConeProperty}). }

 

\subsection{\textcolor{Chapter }{NmzConeDecomposition}}
\logpage{[ 2, 3, 13 ]}\nobreak
\hyperdef{L}{X7B2BD880858784A9}{}
{\noindent\textcolor{FuncColor}{$\triangleright$\enspace\texttt{NmzConeDecomposition({\mdseries\slshape cone})\index{NmzConeDecomposition@\texttt{NmzConeDecomposition}}
\label{NmzConeDecomposition}
}\hfill{\scriptsize (function)}}\\


 This is an alias for \texttt{NmzConeProperty( cone, "ConeDecomposition" );} see \texttt{NmzConeProperty} (\ref{NmzConeProperty}). }

 

\subsection{\textcolor{Chapter }{NmzCongruences}}
\logpage{[ 2, 3, 14 ]}\nobreak
\hyperdef{L}{X7C7798F8847FE06F}{}
{\noindent\textcolor{FuncColor}{$\triangleright$\enspace\texttt{NmzCongruences({\mdseries\slshape cone})\index{NmzCongruences@\texttt{NmzCongruences}}
\label{NmzCongruences}
}\hfill{\scriptsize (function)}}\\
\textbf{\indent Returns:\ }
a matrix whose rows represent the congruences 



 

 The equations, congruences and support hyperplanes together describe the
lattice points of the cone. 

 This is an alias for \texttt{NmzConeProperty( cone, "Congruences" );} see \texttt{NmzConeProperty} (\ref{NmzConeProperty}). }

 

\subsection{\textcolor{Chapter }{NmzCoveringFace}}
\logpage{[ 2, 3, 15 ]}\nobreak
\hyperdef{L}{X7999AC97819BF6EC}{}
{\noindent\textcolor{FuncColor}{$\triangleright$\enspace\texttt{NmzCoveringFace({\mdseries\slshape cone})\index{NmzCoveringFace@\texttt{NmzCoveringFace}}
\label{NmzCoveringFace}
}\hfill{\scriptsize (function)}}\\


 This is an alias for \texttt{NmzConeProperty( cone, "CoveringFace" );} see \texttt{NmzConeProperty} (\ref{NmzConeProperty}). }

 

\subsection{\textcolor{Chapter }{NmzDefaultMode}}
\logpage{[ 2, 3, 16 ]}\nobreak
\hyperdef{L}{X7CE49A02802DE038}{}
{\noindent\textcolor{FuncColor}{$\triangleright$\enspace\texttt{NmzDefaultMode({\mdseries\slshape cone})\index{NmzDefaultMode@\texttt{NmzDefaultMode}}
\label{NmzDefaultMode}
}\hfill{\scriptsize (function)}}\\


 This is an alias for \texttt{NmzConeProperty( cone, "DefaultMode" );} see \texttt{NmzConeProperty} (\ref{NmzConeProperty}). }

 

\subsection{\textcolor{Chapter }{NmzDeg1Elements}}
\logpage{[ 2, 3, 17 ]}\nobreak
\hyperdef{L}{X804F0A437C267569}{}
{\noindent\textcolor{FuncColor}{$\triangleright$\enspace\texttt{NmzDeg1Elements({\mdseries\slshape cone})\index{NmzDeg1Elements@\texttt{NmzDeg1Elements}}
\label{NmzDeg1Elements}
}\hfill{\scriptsize (function)}}\\
\textbf{\indent Returns:\ }
a matrix whose rows are the degree 1 elements 



 

 Requires the presence of a grading. Not allowed in inhomogeneous computations. 

 This is an alias for \texttt{NmzConeProperty( cone, "Deg1Elements" );} see \texttt{NmzConeProperty} (\ref{NmzConeProperty}). }

 

\subsection{\textcolor{Chapter }{NmzDehomogenization}}
\logpage{[ 2, 3, 18 ]}\nobreak
\hyperdef{L}{X81909669839E5AC5}{}
{\noindent\textcolor{FuncColor}{$\triangleright$\enspace\texttt{NmzDehomogenization({\mdseries\slshape cone})\index{NmzDehomogenization@\texttt{NmzDehomogenization}}
\label{NmzDehomogenization}
}\hfill{\scriptsize (function)}}\\
\textbf{\indent Returns:\ }
the dehomgenization vector 



 

 Only for inhomogeneous computations. 

 This is an alias for \texttt{NmzConeProperty( cone, "Dehomogenization" );} see \texttt{NmzConeProperty} (\ref{NmzConeProperty}). }

 

\subsection{\textcolor{Chapter }{NmzDescent}}
\logpage{[ 2, 3, 19 ]}\nobreak
\hyperdef{L}{X818A93CB81759F9B}{}
{\noindent\textcolor{FuncColor}{$\triangleright$\enspace\texttt{NmzDescent({\mdseries\slshape cone})\index{NmzDescent@\texttt{NmzDescent}}
\label{NmzDescent}
}\hfill{\scriptsize (function)}}\\


 This is an alias for \texttt{NmzConeProperty( cone, "Descent" );} see \texttt{NmzConeProperty} (\ref{NmzConeProperty}). }

 

\subsection{\textcolor{Chapter }{NmzDistributedComp}}
\logpage{[ 2, 3, 20 ]}\nobreak
\hyperdef{L}{X85F7F7287C5F7E11}{}
{\noindent\textcolor{FuncColor}{$\triangleright$\enspace\texttt{NmzDistributedComp({\mdseries\slshape cone})\index{NmzDistributedComp@\texttt{NmzDistributedComp}}
\label{NmzDistributedComp}
}\hfill{\scriptsize (function)}}\\


 This is an alias for \texttt{NmzConeProperty( cone, "DistributedComp" );} see \texttt{NmzConeProperty} (\ref{NmzConeProperty}). }

 

\subsection{\textcolor{Chapter }{NmzDualFVector}}
\logpage{[ 2, 3, 21 ]}\nobreak
\hyperdef{L}{X82FF327E7D6109A7}{}
{\noindent\textcolor{FuncColor}{$\triangleright$\enspace\texttt{NmzDualFVector({\mdseries\slshape cone})\index{NmzDualFVector@\texttt{NmzDualFVector}}
\label{NmzDualFVector}
}\hfill{\scriptsize (function)}}\\


 This is an alias for \texttt{NmzConeProperty( cone, "DualFVector" );} see \texttt{NmzConeProperty} (\ref{NmzConeProperty}). }

 

\subsection{\textcolor{Chapter }{NmzDualFaceLattice}}
\logpage{[ 2, 3, 22 ]}\nobreak
\hyperdef{L}{X83E125847D34529B}{}
{\noindent\textcolor{FuncColor}{$\triangleright$\enspace\texttt{NmzDualFaceLattice({\mdseries\slshape cone})\index{NmzDualFaceLattice@\texttt{NmzDualFaceLattice}}
\label{NmzDualFaceLattice}
}\hfill{\scriptsize (function)}}\\


 This is an alias for \texttt{NmzConeProperty( cone, "DualFaceLattice" );} see \texttt{NmzConeProperty} (\ref{NmzConeProperty}). }

 

\subsection{\textcolor{Chapter }{NmzDualIncidence}}
\logpage{[ 2, 3, 23 ]}\nobreak
\hyperdef{L}{X820B11A286A19E8E}{}
{\noindent\textcolor{FuncColor}{$\triangleright$\enspace\texttt{NmzDualIncidence({\mdseries\slshape cone})\index{NmzDualIncidence@\texttt{NmzDualIncidence}}
\label{NmzDualIncidence}
}\hfill{\scriptsize (function)}}\\


 This is an alias for \texttt{NmzConeProperty( cone, "DualIncidence" );} see \texttt{NmzConeProperty} (\ref{NmzConeProperty}). }

 

\subsection{\textcolor{Chapter }{NmzDualMode}}
\logpage{[ 2, 3, 24 ]}\nobreak
\hyperdef{L}{X83D8B1647A5E4C6E}{}
{\noindent\textcolor{FuncColor}{$\triangleright$\enspace\texttt{NmzDualMode({\mdseries\slshape cone})\index{NmzDualMode@\texttt{NmzDualMode}}
\label{NmzDualMode}
}\hfill{\scriptsize (function)}}\\


 This is an alias for \texttt{NmzConeProperty( cone, "DualMode" );} see \texttt{NmzConeProperty} (\ref{NmzConeProperty}). }

 

\subsection{\textcolor{Chapter }{NmzDynamic}}
\logpage{[ 2, 3, 25 ]}\nobreak
\hyperdef{L}{X79C452757E8686FB}{}
{\noindent\textcolor{FuncColor}{$\triangleright$\enspace\texttt{NmzDynamic({\mdseries\slshape cone})\index{NmzDynamic@\texttt{NmzDynamic}}
\label{NmzDynamic}
}\hfill{\scriptsize (function)}}\\


 This is an alias for \texttt{NmzConeProperty( cone, "Dynamic" );} see \texttt{NmzConeProperty} (\ref{NmzConeProperty}). }

 

\subsection{\textcolor{Chapter }{NmzEhrhartQuasiPolynomial}}
\logpage{[ 2, 3, 26 ]}\nobreak
\hyperdef{L}{X7A2C84F0822015AF}{}
{\noindent\textcolor{FuncColor}{$\triangleright$\enspace\texttt{NmzEhrhartQuasiPolynomial({\mdseries\slshape cone})\index{NmzEhrhartQuasiPolynomial@\texttt{NmzEhrhartQuasiPolynomial}}
\label{NmzEhrhartQuasiPolynomial}
}\hfill{\scriptsize (function)}}\\


 This is an alias for \texttt{NmzConeProperty( cone, "EhrhartQuasiPolynomial" );} see \texttt{NmzConeProperty} (\ref{NmzConeProperty}). }

 

\subsection{\textcolor{Chapter }{NmzEhrhartSeries}}
\logpage{[ 2, 3, 27 ]}\nobreak
\hyperdef{L}{X796F86A078BB2E13}{}
{\noindent\textcolor{FuncColor}{$\triangleright$\enspace\texttt{NmzEhrhartSeries({\mdseries\slshape cone})\index{NmzEhrhartSeries@\texttt{NmzEhrhartSeries}}
\label{NmzEhrhartSeries}
}\hfill{\scriptsize (function)}}\\


 

 Supported in Normaliz {\textgreater}= 3.5.0. 

 This is an alias for \texttt{NmzConeProperty( cone, "EhrhartSeries" );} see \texttt{NmzConeProperty} (\ref{NmzConeProperty}). }

 

\subsection{\textcolor{Chapter }{NmzEmbeddingDimension}}
\logpage{[ 2, 3, 28 ]}\nobreak
\hyperdef{L}{X7D96A251809EC941}{}
{\noindent\textcolor{FuncColor}{$\triangleright$\enspace\texttt{NmzEmbeddingDimension({\mdseries\slshape cone})\index{NmzEmbeddingDimension@\texttt{NmzEmbeddingDimension}}
\label{NmzEmbeddingDimension}
}\hfill{\scriptsize (function)}}\\
\textbf{\indent Returns:\ }
the embedding dimension of the cone 



 

 The embedding dimension is the dimension of the space in which the computation
is done. It is the number of components of the output vectors. This value is
always known directly after the creation of the cone. 

 This is an alias for \texttt{NmzConeProperty( cone, "EmbeddingDim" );} see \texttt{NmzConeProperty} (\ref{NmzConeProperty}). }

 

\subsection{\textcolor{Chapter }{NmzEquations}}
\logpage{[ 2, 3, 29 ]}\nobreak
\hyperdef{L}{X7DA58BB08657810D}{}
{\noindent\textcolor{FuncColor}{$\triangleright$\enspace\texttt{NmzEquations({\mdseries\slshape cone})\index{NmzEquations@\texttt{NmzEquations}}
\label{NmzEquations}
}\hfill{\scriptsize (function)}}\\
\textbf{\indent Returns:\ }
a matrix whose rows represent the equations 



 

 The equations cut out the linear space generated by the cone. The equations,
congruences and support hyperplanes together describe the lattice points of
the cone. 

 This is an alias for \texttt{NmzConeProperty( cone, "Equations" );} see \texttt{NmzConeProperty} (\ref{NmzConeProperty}). }

 

\subsection{\textcolor{Chapter }{NmzEuclideanAutomorphisms}}
\logpage{[ 2, 3, 30 ]}\nobreak
\hyperdef{L}{X865B138A80CD79B9}{}
{\noindent\textcolor{FuncColor}{$\triangleright$\enspace\texttt{NmzEuclideanAutomorphisms({\mdseries\slshape cone})\index{NmzEuclideanAutomorphisms@\texttt{NmzEuclideanAutomorphisms}}
\label{NmzEuclideanAutomorphisms}
}\hfill{\scriptsize (function)}}\\


 This is an alias for \texttt{NmzConeProperty( cone, "EuclideanAutomorphisms" );} see \texttt{NmzConeProperty} (\ref{NmzConeProperty}). }

 

\subsection{\textcolor{Chapter }{NmzEuclideanIntegral}}
\logpage{[ 2, 3, 31 ]}\nobreak
\hyperdef{L}{X87919D897817F584}{}
{\noindent\textcolor{FuncColor}{$\triangleright$\enspace\texttt{NmzEuclideanIntegral({\mdseries\slshape cone})\index{NmzEuclideanIntegral@\texttt{NmzEuclideanIntegral}}
\label{NmzEuclideanIntegral}
}\hfill{\scriptsize (function)}}\\


 This is an alias for \texttt{NmzConeProperty( cone, "EuclideanIntegral" );} see \texttt{NmzConeProperty} (\ref{NmzConeProperty}). }

 

\subsection{\textcolor{Chapter }{NmzEuclideanVolume}}
\logpage{[ 2, 3, 32 ]}\nobreak
\hyperdef{L}{X7E55AF3678F4EC46}{}
{\noindent\textcolor{FuncColor}{$\triangleright$\enspace\texttt{NmzEuclideanVolume({\mdseries\slshape cone})\index{NmzEuclideanVolume@\texttt{NmzEuclideanVolume}}
\label{NmzEuclideanVolume}
}\hfill{\scriptsize (function)}}\\


 

 Supported in Normaliz {\textgreater}= 3.5.0. 

 This is an alias for \texttt{NmzConeProperty( cone, "EuclideanVolume" );} see \texttt{NmzConeProperty} (\ref{NmzConeProperty}). }

 

\subsection{\textcolor{Chapter }{NmzExcludedFaces}}
\logpage{[ 2, 3, 33 ]}\nobreak
\hyperdef{L}{X8606AB57813797E5}{}
{\noindent\textcolor{FuncColor}{$\triangleright$\enspace\texttt{NmzExcludedFaces({\mdseries\slshape cone})\index{NmzExcludedFaces@\texttt{NmzExcludedFaces}}
\label{NmzExcludedFaces}
}\hfill{\scriptsize (function)}}\\
\textbf{\indent Returns:\ }
a matrix whose rows represent the excluded faces 



 This is an alias for \texttt{NmzConeProperty( cone, "ExcludedFaces" );} see \texttt{NmzConeProperty} (\ref{NmzConeProperty}). }

 

\subsection{\textcolor{Chapter }{NmzExploitAutomsVectors}}
\logpage{[ 2, 3, 34 ]}\nobreak
\hyperdef{L}{X79E3505A782163F4}{}
{\noindent\textcolor{FuncColor}{$\triangleright$\enspace\texttt{NmzExploitAutomsVectors({\mdseries\slshape cone})\index{NmzExploitAutomsVectors@\texttt{NmzExploitAutomsVectors}}
\label{NmzExploitAutomsVectors}
}\hfill{\scriptsize (function)}}\\


 This is an alias for \texttt{NmzConeProperty( cone, "ExploitAutomsVectors" );} see \texttt{NmzConeProperty} (\ref{NmzConeProperty}). }

 

\subsection{\textcolor{Chapter }{NmzExploitIsosMult}}
\logpage{[ 2, 3, 35 ]}\nobreak
\hyperdef{L}{X862AFE757A5DD829}{}
{\noindent\textcolor{FuncColor}{$\triangleright$\enspace\texttt{NmzExploitIsosMult({\mdseries\slshape cone})\index{NmzExploitIsosMult@\texttt{NmzExploitIsosMult}}
\label{NmzExploitIsosMult}
}\hfill{\scriptsize (function)}}\\


 This is an alias for \texttt{NmzConeProperty( cone, "ExploitIsosMult" );} see \texttt{NmzConeProperty} (\ref{NmzConeProperty}). }

 

\subsection{\textcolor{Chapter }{NmzExternalIndex}}
\logpage{[ 2, 3, 36 ]}\nobreak
\hyperdef{L}{X7BEA1D007DE03002}{}
{\noindent\textcolor{FuncColor}{$\triangleright$\enspace\texttt{NmzExternalIndex({\mdseries\slshape cone})\index{NmzExternalIndex@\texttt{NmzExternalIndex}}
\label{NmzExternalIndex}
}\hfill{\scriptsize (function)}}\\


 This is an alias for \texttt{NmzConeProperty( cone, "ExternalIndex" );} see \texttt{NmzConeProperty} (\ref{NmzConeProperty}). }

 

\subsection{\textcolor{Chapter }{NmzExtremeRays}}
\logpage{[ 2, 3, 37 ]}\nobreak
\hyperdef{L}{X855CB6B5820F28FC}{}
{\noindent\textcolor{FuncColor}{$\triangleright$\enspace\texttt{NmzExtremeRays({\mdseries\slshape cone})\index{NmzExtremeRays@\texttt{NmzExtremeRays}}
\label{NmzExtremeRays}
}\hfill{\scriptsize (function)}}\\
\textbf{\indent Returns:\ }
a matrix whose rows are the extreme rays 



 This is an alias for \texttt{NmzConeProperty( cone, "ExtremeRays" );} see \texttt{NmzConeProperty} (\ref{NmzConeProperty}). }

 

\subsection{\textcolor{Chapter }{NmzExtremeRaysFloat}}
\logpage{[ 2, 3, 38 ]}\nobreak
\hyperdef{L}{X7FB459597A3A1787}{}
{\noindent\textcolor{FuncColor}{$\triangleright$\enspace\texttt{NmzExtremeRaysFloat({\mdseries\slshape cone})\index{NmzExtremeRaysFloat@\texttt{NmzExtremeRaysFloat}}
\label{NmzExtremeRaysFloat}
}\hfill{\scriptsize (function)}}\\


 This is an alias for \texttt{NmzConeProperty( cone, "ExtremeRaysFloat" );} see \texttt{NmzConeProperty} (\ref{NmzConeProperty}). }

 

\subsection{\textcolor{Chapter }{NmzFVector}}
\logpage{[ 2, 3, 39 ]}\nobreak
\hyperdef{L}{X8084862B7E5EB685}{}
{\noindent\textcolor{FuncColor}{$\triangleright$\enspace\texttt{NmzFVector({\mdseries\slshape cone})\index{NmzFVector@\texttt{NmzFVector}}
\label{NmzFVector}
}\hfill{\scriptsize (function)}}\\


 

 Supported in Normaliz {\textgreater}= 3.7.0. 

 This is an alias for \texttt{NmzConeProperty( cone, "FVector" );} see \texttt{NmzConeProperty} (\ref{NmzConeProperty}). }

 

\subsection{\textcolor{Chapter }{NmzFaceLattice}}
\logpage{[ 2, 3, 40 ]}\nobreak
\hyperdef{L}{X87916CA882C9B4CC}{}
{\noindent\textcolor{FuncColor}{$\triangleright$\enspace\texttt{NmzFaceLattice({\mdseries\slshape cone})\index{NmzFaceLattice@\texttt{NmzFaceLattice}}
\label{NmzFaceLattice}
}\hfill{\scriptsize (function)}}\\


 

 Supported in Normaliz {\textgreater}= 3.7.0. 

 This is an alias for \texttt{NmzConeProperty( cone, "FaceLattice" );} see \texttt{NmzConeProperty} (\ref{NmzConeProperty}). }

 

\subsection{\textcolor{Chapter }{NmzFixedPrecision}}
\logpage{[ 2, 3, 41 ]}\nobreak
\hyperdef{L}{X836BB3CE7C97B141}{}
{\noindent\textcolor{FuncColor}{$\triangleright$\enspace\texttt{NmzFixedPrecision({\mdseries\slshape cone})\index{NmzFixedPrecision@\texttt{NmzFixedPrecision}}
\label{NmzFixedPrecision}
}\hfill{\scriptsize (function)}}\\


 This is an alias for \texttt{NmzConeProperty( cone, "FixedPrecision" );} see \texttt{NmzConeProperty} (\ref{NmzConeProperty}). }

 

\subsection{\textcolor{Chapter }{NmzFullConeDynamic}}
\logpage{[ 2, 3, 42 ]}\nobreak
\hyperdef{L}{X7CA11F1787307478}{}
{\noindent\textcolor{FuncColor}{$\triangleright$\enspace\texttt{NmzFullConeDynamic({\mdseries\slshape cone})\index{NmzFullConeDynamic@\texttt{NmzFullConeDynamic}}
\label{NmzFullConeDynamic}
}\hfill{\scriptsize (function)}}\\


 This is an alias for \texttt{NmzConeProperty( cone, "FullConeDynamic" );} see \texttt{NmzConeProperty} (\ref{NmzConeProperty}). }

 

\subsection{\textcolor{Chapter }{NmzGeneratorOfInterior}}
\logpage{[ 2, 3, 43 ]}\nobreak
\hyperdef{L}{X8782055C7F394B20}{}
{\noindent\textcolor{FuncColor}{$\triangleright$\enspace\texttt{NmzGeneratorOfInterior({\mdseries\slshape cone})\index{NmzGeneratorOfInterior@\texttt{NmzGeneratorOfInterior}}
\label{NmzGeneratorOfInterior}
}\hfill{\scriptsize (function)}}\\
\textbf{\indent Returns:\ }
a vector representing the generator of the interior of \mbox{\texttt{\mdseries\slshape cone}} 



 

 If \mbox{\texttt{\mdseries\slshape cone}} is Gorenstein, this function returns the generator of the interior of \mbox{\texttt{\mdseries\slshape cone}}. If \mbox{\texttt{\mdseries\slshape cone}} is not Gorenstein, an error is raised. 

 This is an alias for \texttt{NmzConeProperty( cone, "GeneratorOfInterior" );} see \texttt{NmzConeProperty} (\ref{NmzConeProperty}). }

 

\subsection{\textcolor{Chapter }{NmzGenerators}}
\logpage{[ 2, 3, 44 ]}\nobreak
\hyperdef{L}{X87E77F007E199FCA}{}
{\noindent\textcolor{FuncColor}{$\triangleright$\enspace\texttt{NmzGenerators({\mdseries\slshape cone})\index{NmzGenerators@\texttt{NmzGenerators}}
\label{NmzGenerators}
}\hfill{\scriptsize (function)}}\\
\textbf{\indent Returns:\ }
a matrix whose rows are the generators of \mbox{\texttt{\mdseries\slshape cone}} 



 This is an alias for \texttt{NmzConeProperty( cone, "Generators" );} see \texttt{NmzConeProperty} (\ref{NmzConeProperty}). }

 

\subsection{\textcolor{Chapter }{NmzGrading}}
\logpage{[ 2, 3, 45 ]}\nobreak
\hyperdef{L}{X83DC150E7C7A3AF7}{}
{\noindent\textcolor{FuncColor}{$\triangleright$\enspace\texttt{NmzGrading({\mdseries\slshape cone})\index{NmzGrading@\texttt{NmzGrading}}
\label{NmzGrading}
}\hfill{\scriptsize (function)}}\\
\textbf{\indent Returns:\ }
the grading vector 



 This is an alias for \texttt{NmzConeProperty( cone, "Grading" );} see \texttt{NmzConeProperty} (\ref{NmzConeProperty}). }

 

\subsection{\textcolor{Chapter }{NmzGradingDenom}}
\logpage{[ 2, 3, 46 ]}\nobreak
\hyperdef{L}{X812A9E21823DCC42}{}
{\noindent\textcolor{FuncColor}{$\triangleright$\enspace\texttt{NmzGradingDenom({\mdseries\slshape cone})\index{NmzGradingDenom@\texttt{NmzGradingDenom}}
\label{NmzGradingDenom}
}\hfill{\scriptsize (function)}}\\


 This is an alias for \texttt{NmzConeProperty( cone, "GradingDenom" );} see \texttt{NmzConeProperty} (\ref{NmzConeProperty}). }

 

\subsection{\textcolor{Chapter }{NmzGradingIsPositive}}
\logpage{[ 2, 3, 47 ]}\nobreak
\hyperdef{L}{X804F6DF27FA1FD5E}{}
{\noindent\textcolor{FuncColor}{$\triangleright$\enspace\texttt{NmzGradingIsPositive({\mdseries\slshape cone})\index{NmzGradingIsPositive@\texttt{NmzGradingIsPositive}}
\label{NmzGradingIsPositive}
}\hfill{\scriptsize (function)}}\\


 This is an alias for \texttt{NmzConeProperty( cone, "GradingIsPositive" );} see \texttt{NmzConeProperty} (\ref{NmzConeProperty}). }

 

\subsection{\textcolor{Chapter }{NmzHSOP}}
\logpage{[ 2, 3, 48 ]}\nobreak
\hyperdef{L}{X86D252817DDF4689}{}
{\noindent\textcolor{FuncColor}{$\triangleright$\enspace\texttt{NmzHSOP({\mdseries\slshape cone})\index{NmzHSOP@\texttt{NmzHSOP}}
\label{NmzHSOP}
}\hfill{\scriptsize (function)}}\\


 This is an alias for \texttt{NmzConeProperty( cone, "HSOP" );} see \texttt{NmzConeProperty} (\ref{NmzConeProperty}). }

 

\subsection{\textcolor{Chapter }{NmzHilbertBasis}}
\logpage{[ 2, 3, 49 ]}\nobreak
\hyperdef{L}{X7A11497282B4831D}{}
{\noindent\textcolor{FuncColor}{$\triangleright$\enspace\texttt{NmzHilbertBasis({\mdseries\slshape cone})\index{NmzHilbertBasis@\texttt{NmzHilbertBasis}}
\label{NmzHilbertBasis}
}\hfill{\scriptsize (function)}}\\
\textbf{\indent Returns:\ }
a matrix whose rows are the Hilbert basis elements 



 This is an alias for \texttt{NmzConeProperty( cone, "HilbertBasis" );} see \texttt{NmzConeProperty} (\ref{NmzConeProperty}). }

 

\subsection{\textcolor{Chapter }{NmzHilbertQuasiPolynomial}}
\logpage{[ 2, 3, 50 ]}\nobreak
\hyperdef{L}{X8405B3167F4ED99C}{}
{\noindent\textcolor{FuncColor}{$\triangleright$\enspace\texttt{NmzHilbertQuasiPolynomial({\mdseries\slshape cone})\index{NmzHilbertQuasiPolynomial@\texttt{NmzHilbertQuasiPolynomial}}
\label{NmzHilbertQuasiPolynomial}
}\hfill{\scriptsize (function)}}\\
\textbf{\indent Returns:\ }
the Hilbert function as a quasipolynomial 



 

 The Hilbert function counts the lattice points degreewise. The result is a
quasipolynomial $Q$, that is, a polynomial with periodic coefficients. It is given as list of
polynomials $P_0, \ldots, P_{(p-1)}$ such that $Q(i) = P_{(i \bmod p)} (i)$. 

 This is an alias for \texttt{NmzConeProperty( cone, "HilbertQuasiPolynomial" );} see \texttt{NmzConeProperty} (\ref{NmzConeProperty}). }

 

\subsection{\textcolor{Chapter }{NmzHilbertSeries}}
\logpage{[ 2, 3, 51 ]}\nobreak
\hyperdef{L}{X7FC3F68185D5E220}{}
{\noindent\textcolor{FuncColor}{$\triangleright$\enspace\texttt{NmzHilbertSeries({\mdseries\slshape cone})\index{NmzHilbertSeries@\texttt{NmzHilbertSeries}}
\label{NmzHilbertSeries}
}\hfill{\scriptsize (function)}}\\
\textbf{\indent Returns:\ }
the Hilbert series as rational function 



 

 The result consists of a list with two entries. The first is the numerator
polynomial. In inhomogeneous computations this can also be a Laurent
polynomial. The second list entry represents the denominator. It is a list of
pairs $[k_i, l_i]$. Such a pair represents the factor $(1-t^{k_i})^{l_i}$. 

 This is an alias for \texttt{NmzConeProperty( cone, "HilbertSeries" );} see \texttt{NmzConeProperty} (\ref{NmzConeProperty}). }

 

\subsection{\textcolor{Chapter }{NmzIncidence}}
\logpage{[ 2, 3, 52 ]}\nobreak
\hyperdef{L}{X81B24C9187682ACC}{}
{\noindent\textcolor{FuncColor}{$\triangleright$\enspace\texttt{NmzIncidence({\mdseries\slshape cone})\index{NmzIncidence@\texttt{NmzIncidence}}
\label{NmzIncidence}
}\hfill{\scriptsize (function)}}\\


 

 Supported in Normaliz {\textgreater}= 3.8.0. 

 This is an alias for \texttt{NmzConeProperty( cone, "Incidence" );} see \texttt{NmzConeProperty} (\ref{NmzConeProperty}). }

 

\subsection{\textcolor{Chapter }{NmzInclusionExclusionData}}
\logpage{[ 2, 3, 53 ]}\nobreak
\hyperdef{L}{X8491B203791F0526}{}
{\noindent\textcolor{FuncColor}{$\triangleright$\enspace\texttt{NmzInclusionExclusionData({\mdseries\slshape cone})\index{NmzInclusionExclusionData@\texttt{NmzInclusionExclusionData}}
\label{NmzInclusionExclusionData}
}\hfill{\scriptsize (function)}}\\
\textbf{\indent Returns:\ }
inclusion\texttt{\symbol{45}}exclusion data 



 

 List of faces which are internally have been used in the
inclusion\texttt{\symbol{45}}exclusion scheme. Given as a list pairs. The
first pair entry is a key of generators contained in the face (compare also \texttt{NmzTriangulation} (\ref{NmzTriangulation})) and the multiplicity with which it was considered. Only available with
excluded faces or strict constraints as input. 

 This is an alias for \texttt{NmzConeProperty( cone, "InclusionExclusionData" );} see \texttt{NmzConeProperty} (\ref{NmzConeProperty}). }

 

\subsection{\textcolor{Chapter }{NmzInputAutomorphisms}}
\logpage{[ 2, 3, 54 ]}\nobreak
\hyperdef{L}{X7C6C9DAE843CC8B4}{}
{\noindent\textcolor{FuncColor}{$\triangleright$\enspace\texttt{NmzInputAutomorphisms({\mdseries\slshape cone})\index{NmzInputAutomorphisms@\texttt{NmzInputAutomorphisms}}
\label{NmzInputAutomorphisms}
}\hfill{\scriptsize (function)}}\\


 This is an alias for \texttt{NmzConeProperty( cone, "InputAutomorphisms" );} see \texttt{NmzConeProperty} (\ref{NmzConeProperty}). }

 

\subsection{\textcolor{Chapter }{NmzIntegerHull}}
\logpage{[ 2, 3, 55 ]}\nobreak
\hyperdef{L}{X83425BF67C13B187}{}
{\noindent\textcolor{FuncColor}{$\triangleright$\enspace\texttt{NmzIntegerHull({\mdseries\slshape cone})\index{NmzIntegerHull@\texttt{NmzIntegerHull}}
\label{NmzIntegerHull}
}\hfill{\scriptsize (function)}}\\


 This is an alias for \texttt{NmzConeProperty( cone, "IntegerHull" );} see \texttt{NmzConeProperty} (\ref{NmzConeProperty}). }

 

\subsection{\textcolor{Chapter }{NmzIntegral}}
\logpage{[ 2, 3, 56 ]}\nobreak
\hyperdef{L}{X7E2B4A657D716227}{}
{\noindent\textcolor{FuncColor}{$\triangleright$\enspace\texttt{NmzIntegral({\mdseries\slshape cone})\index{NmzIntegral@\texttt{NmzIntegral}}
\label{NmzIntegral}
}\hfill{\scriptsize (function)}}\\


 This is an alias for \texttt{NmzConeProperty( cone, "Integral" );} see \texttt{NmzConeProperty} (\ref{NmzConeProperty}). }

 

\subsection{\textcolor{Chapter }{NmzInternalIndex}}
\logpage{[ 2, 3, 57 ]}\nobreak
\hyperdef{L}{X7C73E21E87B75DDB}{}
{\noindent\textcolor{FuncColor}{$\triangleright$\enspace\texttt{NmzInternalIndex({\mdseries\slshape cone})\index{NmzInternalIndex@\texttt{NmzInternalIndex}}
\label{NmzInternalIndex}
}\hfill{\scriptsize (function)}}\\


 This is an alias for \texttt{NmzConeProperty( cone, "InternalIndex" );} see \texttt{NmzConeProperty} (\ref{NmzConeProperty}). }

 

\subsection{\textcolor{Chapter }{NmzIsDeg1ExtremeRays}}
\logpage{[ 2, 3, 58 ]}\nobreak
\hyperdef{L}{X847176FE85D292D5}{}
{\noindent\textcolor{FuncColor}{$\triangleright$\enspace\texttt{NmzIsDeg1ExtremeRays({\mdseries\slshape cone})\index{NmzIsDeg1ExtremeRays@\texttt{NmzIsDeg1ExtremeRays}}
\label{NmzIsDeg1ExtremeRays}
}\hfill{\scriptsize (function)}}\\
\textbf{\indent Returns:\ }
\texttt{true} if all extreme rays have degree 1; \texttt{false} otherwise 



 This is an alias for \texttt{NmzConeProperty( cone, "IsDeg1ExtremeRays" );} see \texttt{NmzConeProperty} (\ref{NmzConeProperty}). }

 

\subsection{\textcolor{Chapter }{NmzIsDeg1HilbertBasis}}
\logpage{[ 2, 3, 59 ]}\nobreak
\hyperdef{L}{X7846999882E09FBC}{}
{\noindent\textcolor{FuncColor}{$\triangleright$\enspace\texttt{NmzIsDeg1HilbertBasis({\mdseries\slshape cone})\index{NmzIsDeg1HilbertBasis@\texttt{NmzIsDeg1HilbertBasis}}
\label{NmzIsDeg1HilbertBasis}
}\hfill{\scriptsize (function)}}\\
\textbf{\indent Returns:\ }
\texttt{true} if all Hilbert basis elements have degree 1; \texttt{false} otherwise 



 This is an alias for \texttt{NmzConeProperty( cone, "IsDeg1HilbertBasis" );} see \texttt{NmzConeProperty} (\ref{NmzConeProperty}). }

 

\subsection{\textcolor{Chapter }{NmzIsEmptySemiOpen}}
\logpage{[ 2, 3, 60 ]}\nobreak
\hyperdef{L}{X84A71139794B76F9}{}
{\noindent\textcolor{FuncColor}{$\triangleright$\enspace\texttt{NmzIsEmptySemiOpen({\mdseries\slshape cone})\index{NmzIsEmptySemiOpen@\texttt{NmzIsEmptySemiOpen}}
\label{NmzIsEmptySemiOpen}
}\hfill{\scriptsize (function)}}\\


 This is an alias for \texttt{NmzConeProperty( cone, "IsEmptySemiOpen" );} see \texttt{NmzConeProperty} (\ref{NmzConeProperty}). }

 

\subsection{\textcolor{Chapter }{NmzIsGorenstein}}
\logpage{[ 2, 3, 61 ]}\nobreak
\hyperdef{L}{X848E74AA86BAA5B9}{}
{\noindent\textcolor{FuncColor}{$\triangleright$\enspace\texttt{NmzIsGorenstein({\mdseries\slshape cone})\index{NmzIsGorenstein@\texttt{NmzIsGorenstein}}
\label{NmzIsGorenstein}
}\hfill{\scriptsize (function)}}\\
\textbf{\indent Returns:\ }
whether the cone is Gorenstein 



 

 Returns true if \mbox{\texttt{\mdseries\slshape cone}} is Gorenstein, false otherwise. 

 This is an alias for \texttt{NmzConeProperty( cone, "IsGorenstein" );} see \texttt{NmzConeProperty} (\ref{NmzConeProperty}). }

 

\subsection{\textcolor{Chapter }{NmzIsInhomogeneous}}
\logpage{[ 2, 3, 62 ]}\nobreak
\hyperdef{L}{X82EF925E7CF91FDD}{}
{\noindent\textcolor{FuncColor}{$\triangleright$\enspace\texttt{NmzIsInhomogeneous({\mdseries\slshape cone})\index{NmzIsInhomogeneous@\texttt{NmzIsInhomogeneous}}
\label{NmzIsInhomogeneous}
}\hfill{\scriptsize (function)}}\\
\textbf{\indent Returns:\ }
whether the cone is inhomogeneous 



 

 This value is always known directly after the creation of the cone. 

 This is an alias for \texttt{NmzConeProperty( cone, "IsInhomogeneous" );} see \texttt{NmzConeProperty} (\ref{NmzConeProperty}). }

 

\subsection{\textcolor{Chapter }{NmzIsIntegrallyClosed}}
\logpage{[ 2, 3, 63 ]}\nobreak
\hyperdef{L}{X8315ACBC797CF531}{}
{\noindent\textcolor{FuncColor}{$\triangleright$\enspace\texttt{NmzIsIntegrallyClosed({\mdseries\slshape cone})\index{NmzIsIntegrallyClosed@\texttt{NmzIsIntegrallyClosed}}
\label{NmzIsIntegrallyClosed}
}\hfill{\scriptsize (function)}}\\
\textbf{\indent Returns:\ }
\texttt{true} if the cone is integrally closed; \texttt{false} otherwise 



 

 It is integrally closed when the Hilbert basis is a subset of the original
monoid generators. So it is only computable if we have original monoid
generators. 

 This is an alias for \texttt{NmzConeProperty( cone, "IsIntegrallyClosed" );} see \texttt{NmzConeProperty} (\ref{NmzConeProperty}). }

 

\subsection{\textcolor{Chapter }{NmzIsPointed}}
\logpage{[ 2, 3, 64 ]}\nobreak
\hyperdef{L}{X863887AA852ACA6F}{}
{\noindent\textcolor{FuncColor}{$\triangleright$\enspace\texttt{NmzIsPointed({\mdseries\slshape cone})\index{NmzIsPointed@\texttt{NmzIsPointed}}
\label{NmzIsPointed}
}\hfill{\scriptsize (function)}}\\
\textbf{\indent Returns:\ }
\texttt{true} if the cone is pointed; \texttt{false} otherwise 



 This is an alias for \texttt{NmzConeProperty( cone, "IsPointed" );} see \texttt{NmzConeProperty} (\ref{NmzConeProperty}). }

 

\subsection{\textcolor{Chapter }{NmzIsReesPrimary}}
\logpage{[ 2, 3, 65 ]}\nobreak
\hyperdef{L}{X80CAC21686223F6F}{}
{\noindent\textcolor{FuncColor}{$\triangleright$\enspace\texttt{NmzIsReesPrimary({\mdseries\slshape cone})\index{NmzIsReesPrimary@\texttt{NmzIsReesPrimary}}
\label{NmzIsReesPrimary}
}\hfill{\scriptsize (function)}}\\
\textbf{\indent Returns:\ }
\texttt{true} if is the monomial ideal is primary to the irrelevant maximal ideal, \texttt{false} otherwise 



 

 Only used with the input type \texttt{rees{\textunderscore}algebra}. 

 This is an alias for \texttt{NmzConeProperty( cone, "IsReesPrimary" );} see \texttt{NmzConeProperty} (\ref{NmzConeProperty}). }

 

\subsection{\textcolor{Chapter }{NmzIsTriangulationNested}}
\logpage{[ 2, 3, 66 ]}\nobreak
\hyperdef{L}{X7EEE4D2F86463F21}{}
{\noindent\textcolor{FuncColor}{$\triangleright$\enspace\texttt{NmzIsTriangulationNested({\mdseries\slshape cone})\index{NmzIsTriangulationNested@\texttt{NmzIsTriangulationNested}}
\label{NmzIsTriangulationNested}
}\hfill{\scriptsize (function)}}\\


 This is an alias for \texttt{NmzConeProperty( cone, "IsTriangulationNested" );} see \texttt{NmzConeProperty} (\ref{NmzConeProperty}). }

 

\subsection{\textcolor{Chapter }{NmzIsTriangulationPartial}}
\logpage{[ 2, 3, 67 ]}\nobreak
\hyperdef{L}{X8159CF6885A164E1}{}
{\noindent\textcolor{FuncColor}{$\triangleright$\enspace\texttt{NmzIsTriangulationPartial({\mdseries\slshape cone})\index{NmzIsTriangulationPartial@\texttt{NmzIsTriangulationPartial}}
\label{NmzIsTriangulationPartial}
}\hfill{\scriptsize (function)}}\\


 This is an alias for \texttt{NmzConeProperty( cone, "IsTriangulationPartial" );} see \texttt{NmzConeProperty} (\ref{NmzConeProperty}). }

 

\subsection{\textcolor{Chapter }{NmzKeepOrder}}
\logpage{[ 2, 3, 68 ]}\nobreak
\hyperdef{L}{X84DB162079D98024}{}
{\noindent\textcolor{FuncColor}{$\triangleright$\enspace\texttt{NmzKeepOrder({\mdseries\slshape cone})\index{NmzKeepOrder@\texttt{NmzKeepOrder}}
\label{NmzKeepOrder}
}\hfill{\scriptsize (function)}}\\


 This is an alias for \texttt{NmzConeProperty( cone, "KeepOrder" );} see \texttt{NmzConeProperty} (\ref{NmzConeProperty}). }

 

\subsection{\textcolor{Chapter }{NmzLatticePointTriangulation}}
\logpage{[ 2, 3, 69 ]}\nobreak
\hyperdef{L}{X870368D0819E8170}{}
{\noindent\textcolor{FuncColor}{$\triangleright$\enspace\texttt{NmzLatticePointTriangulation({\mdseries\slshape cone})\index{NmzLatticePointTriangulation@\texttt{NmzLatticePointTriangulation}}
\label{NmzLatticePointTriangulation}
}\hfill{\scriptsize (function)}}\\


 This is an alias for \texttt{NmzConeProperty( cone, "LatticePointTriangulation" );} see \texttt{NmzConeProperty} (\ref{NmzConeProperty}). }

 

\subsection{\textcolor{Chapter }{NmzLatticePoints}}
\logpage{[ 2, 3, 70 ]}\nobreak
\hyperdef{L}{X78E297A8795F0074}{}
{\noindent\textcolor{FuncColor}{$\triangleright$\enspace\texttt{NmzLatticePoints({\mdseries\slshape cone})\index{NmzLatticePoints@\texttt{NmzLatticePoints}}
\label{NmzLatticePoints}
}\hfill{\scriptsize (function)}}\\


 This is an alias for \texttt{NmzConeProperty( cone, "LatticePoints" );} see \texttt{NmzConeProperty} (\ref{NmzConeProperty}). }

 

\subsection{\textcolor{Chapter }{NmzMaximalSubspace}}
\logpage{[ 2, 3, 71 ]}\nobreak
\hyperdef{L}{X79D80E6081737899}{}
{\noindent\textcolor{FuncColor}{$\triangleright$\enspace\texttt{NmzMaximalSubspace({\mdseries\slshape cone})\index{NmzMaximalSubspace@\texttt{NmzMaximalSubspace}}
\label{NmzMaximalSubspace}
}\hfill{\scriptsize (function)}}\\
\textbf{\indent Returns:\ }
a matrix whose rows generate the maximale linear subspace 



 This is an alias for \texttt{NmzConeProperty( cone, "MaximalSubspace" );} see \texttt{NmzConeProperty} (\ref{NmzConeProperty}). }

 

\subsection{\textcolor{Chapter }{NmzModuleGenerators}}
\logpage{[ 2, 3, 72 ]}\nobreak
\hyperdef{L}{X7CDCEC40814035A8}{}
{\noindent\textcolor{FuncColor}{$\triangleright$\enspace\texttt{NmzModuleGenerators({\mdseries\slshape cone})\index{NmzModuleGenerators@\texttt{NmzModuleGenerators}}
\label{NmzModuleGenerators}
}\hfill{\scriptsize (function)}}\\
\textbf{\indent Returns:\ }
a matrix whose rows are the module generators 



 This is an alias for \texttt{NmzConeProperty( cone, "ModuleGenerators" );} see \texttt{NmzConeProperty} (\ref{NmzConeProperty}). }

 

\subsection{\textcolor{Chapter }{NmzModuleGeneratorsOverOriginalMonoid}}
\logpage{[ 2, 3, 73 ]}\nobreak
\hyperdef{L}{X871D59BC83D7CADA}{}
{\noindent\textcolor{FuncColor}{$\triangleright$\enspace\texttt{NmzModuleGeneratorsOverOriginalMonoid({\mdseries\slshape cone})\index{NmzModuleGeneratorsOverOriginalMonoid@\texttt{Nmz}\-\texttt{Module}\-\texttt{Generators}\-\texttt{Over}\-\texttt{Original}\-\texttt{Monoid}}
\label{NmzModuleGeneratorsOverOriginalMonoid}
}\hfill{\scriptsize (function)}}\\
\textbf{\indent Returns:\ }
a matrix whose rows are the module generators over the original monoid 



 

 A minimal system of generators of the integral closure over the original
monoid. Requires the existence of original monoid generators. Not allowed in
inhomogeneous computations. 

 This is an alias for \texttt{NmzConeProperty( cone, "ModuleGeneratorsOverOriginalMonoid" );} see \texttt{NmzConeProperty} (\ref{NmzConeProperty}). }

 

\subsection{\textcolor{Chapter }{NmzModuleRank}}
\logpage{[ 2, 3, 74 ]}\nobreak
\hyperdef{L}{X81CD718A78E72475}{}
{\noindent\textcolor{FuncColor}{$\triangleright$\enspace\texttt{NmzModuleRank({\mdseries\slshape cone})\index{NmzModuleRank@\texttt{NmzModuleRank}}
\label{NmzModuleRank}
}\hfill{\scriptsize (function)}}\\
\textbf{\indent Returns:\ }
the rank of the module of lattice points in the polyhedron as a module over
the recession monoid 



 

 Only for inhomogeneous computations. 

 This is an alias for \texttt{NmzConeProperty( cone, "ModuleRank" );} see \texttt{NmzConeProperty} (\ref{NmzConeProperty}). }

 

\subsection{\textcolor{Chapter }{NmzMultiplicity}}
\logpage{[ 2, 3, 75 ]}\nobreak
\hyperdef{L}{X7BEE5536807AD064}{}
{\noindent\textcolor{FuncColor}{$\triangleright$\enspace\texttt{NmzMultiplicity({\mdseries\slshape cone})\index{NmzMultiplicity@\texttt{NmzMultiplicity}}
\label{NmzMultiplicity}
}\hfill{\scriptsize (function)}}\\


 This is an alias for \texttt{NmzConeProperty( cone, "Multiplicity" );} see \texttt{NmzConeProperty} (\ref{NmzConeProperty}). }

 

\subsection{\textcolor{Chapter }{NmzNoBottomDec}}
\logpage{[ 2, 3, 76 ]}\nobreak
\hyperdef{L}{X7D14D7D17BC476ED}{}
{\noindent\textcolor{FuncColor}{$\triangleright$\enspace\texttt{NmzNoBottomDec({\mdseries\slshape cone})\index{NmzNoBottomDec@\texttt{NmzNoBottomDec}}
\label{NmzNoBottomDec}
}\hfill{\scriptsize (function)}}\\


 This is an alias for \texttt{NmzConeProperty( cone, "NoBottomDec" );} see \texttt{NmzConeProperty} (\ref{NmzConeProperty}). }

 

\subsection{\textcolor{Chapter }{NmzNoDescent}}
\logpage{[ 2, 3, 77 ]}\nobreak
\hyperdef{L}{X8256AFFA7AD3DB1D}{}
{\noindent\textcolor{FuncColor}{$\triangleright$\enspace\texttt{NmzNoDescent({\mdseries\slshape cone})\index{NmzNoDescent@\texttt{NmzNoDescent}}
\label{NmzNoDescent}
}\hfill{\scriptsize (function)}}\\


 This is an alias for \texttt{NmzConeProperty( cone, "NoDescent" );} see \texttt{NmzConeProperty} (\ref{NmzConeProperty}). }

 

\subsection{\textcolor{Chapter }{NmzNoGradingDenom}}
\logpage{[ 2, 3, 78 ]}\nobreak
\hyperdef{L}{X7E2A036A78414435}{}
{\noindent\textcolor{FuncColor}{$\triangleright$\enspace\texttt{NmzNoGradingDenom({\mdseries\slshape cone})\index{NmzNoGradingDenom@\texttt{NmzNoGradingDenom}}
\label{NmzNoGradingDenom}
}\hfill{\scriptsize (function)}}\\


 This is an alias for \texttt{NmzConeProperty( cone, "NoGradingDenom" );} see \texttt{NmzConeProperty} (\ref{NmzConeProperty}). }

 

\subsection{\textcolor{Chapter }{NmzNoLLL}}
\logpage{[ 2, 3, 79 ]}\nobreak
\hyperdef{L}{X7B6047317EDBB5F2}{}
{\noindent\textcolor{FuncColor}{$\triangleright$\enspace\texttt{NmzNoLLL({\mdseries\slshape cone})\index{NmzNoLLL@\texttt{NmzNoLLL}}
\label{NmzNoLLL}
}\hfill{\scriptsize (function)}}\\


 This is an alias for \texttt{NmzConeProperty( cone, "NoLLL" );} see \texttt{NmzConeProperty} (\ref{NmzConeProperty}). }

 

\subsection{\textcolor{Chapter }{NmzNoNestedTri}}
\logpage{[ 2, 3, 80 ]}\nobreak
\hyperdef{L}{X7ADD635981E04C15}{}
{\noindent\textcolor{FuncColor}{$\triangleright$\enspace\texttt{NmzNoNestedTri({\mdseries\slshape cone})\index{NmzNoNestedTri@\texttt{NmzNoNestedTri}}
\label{NmzNoNestedTri}
}\hfill{\scriptsize (function)}}\\


 This is an alias for \texttt{NmzConeProperty( cone, "NoNestedTri" );} see \texttt{NmzConeProperty} (\ref{NmzConeProperty}). }

 

\subsection{\textcolor{Chapter }{NmzNoPeriodBound}}
\logpage{[ 2, 3, 81 ]}\nobreak
\hyperdef{L}{X802E5F2C84870D9A}{}
{\noindent\textcolor{FuncColor}{$\triangleright$\enspace\texttt{NmzNoPeriodBound({\mdseries\slshape cone})\index{NmzNoPeriodBound@\texttt{NmzNoPeriodBound}}
\label{NmzNoPeriodBound}
}\hfill{\scriptsize (function)}}\\


 This is an alias for \texttt{NmzConeProperty( cone, "NoPeriodBound" );} see \texttt{NmzConeProperty} (\ref{NmzConeProperty}). }

 

\subsection{\textcolor{Chapter }{NmzNoProjection}}
\logpage{[ 2, 3, 82 ]}\nobreak
\hyperdef{L}{X844FFA72847DEF0D}{}
{\noindent\textcolor{FuncColor}{$\triangleright$\enspace\texttt{NmzNoProjection({\mdseries\slshape cone})\index{NmzNoProjection@\texttt{NmzNoProjection}}
\label{NmzNoProjection}
}\hfill{\scriptsize (function)}}\\


 This is an alias for \texttt{NmzConeProperty( cone, "NoProjection" );} see \texttt{NmzConeProperty} (\ref{NmzConeProperty}). }

 

\subsection{\textcolor{Chapter }{NmzNoRelax}}
\logpage{[ 2, 3, 83 ]}\nobreak
\hyperdef{L}{X784B65D984012D39}{}
{\noindent\textcolor{FuncColor}{$\triangleright$\enspace\texttt{NmzNoRelax({\mdseries\slshape cone})\index{NmzNoRelax@\texttt{NmzNoRelax}}
\label{NmzNoRelax}
}\hfill{\scriptsize (function)}}\\


 This is an alias for \texttt{NmzConeProperty( cone, "NoRelax" );} see \texttt{NmzConeProperty} (\ref{NmzConeProperty}). }

 

\subsection{\textcolor{Chapter }{NmzNoSignedDec}}
\logpage{[ 2, 3, 84 ]}\nobreak
\hyperdef{L}{X86F5AD227B0A8DAB}{}
{\noindent\textcolor{FuncColor}{$\triangleright$\enspace\texttt{NmzNoSignedDec({\mdseries\slshape cone})\index{NmzNoSignedDec@\texttt{NmzNoSignedDec}}
\label{NmzNoSignedDec}
}\hfill{\scriptsize (function)}}\\


 This is an alias for \texttt{NmzConeProperty( cone, "NoSignedDec" );} see \texttt{NmzConeProperty} (\ref{NmzConeProperty}). }

 

\subsection{\textcolor{Chapter }{NmzNoSubdivision}}
\logpage{[ 2, 3, 85 ]}\nobreak
\hyperdef{L}{X8274F8F57816A9B5}{}
{\noindent\textcolor{FuncColor}{$\triangleright$\enspace\texttt{NmzNoSubdivision({\mdseries\slshape cone})\index{NmzNoSubdivision@\texttt{NmzNoSubdivision}}
\label{NmzNoSubdivision}
}\hfill{\scriptsize (function)}}\\


 This is an alias for \texttt{NmzConeProperty( cone, "NoSubdivision" );} see \texttt{NmzConeProperty} (\ref{NmzConeProperty}). }

 

\subsection{\textcolor{Chapter }{NmzNoSymmetrization}}
\logpage{[ 2, 3, 86 ]}\nobreak
\hyperdef{L}{X87BEBC4883D9E4CF}{}
{\noindent\textcolor{FuncColor}{$\triangleright$\enspace\texttt{NmzNoSymmetrization({\mdseries\slshape cone})\index{NmzNoSymmetrization@\texttt{NmzNoSymmetrization}}
\label{NmzNoSymmetrization}
}\hfill{\scriptsize (function)}}\\


 This is an alias for \texttt{NmzConeProperty( cone, "NoSymmetrization" );} see \texttt{NmzConeProperty} (\ref{NmzConeProperty}). }

 

\subsection{\textcolor{Chapter }{NmzNumberLatticePoints}}
\logpage{[ 2, 3, 87 ]}\nobreak
\hyperdef{L}{X812E448A7C49CCD3}{}
{\noindent\textcolor{FuncColor}{$\triangleright$\enspace\texttt{NmzNumberLatticePoints({\mdseries\slshape cone})\index{NmzNumberLatticePoints@\texttt{NmzNumberLatticePoints}}
\label{NmzNumberLatticePoints}
}\hfill{\scriptsize (function)}}\\


 This is an alias for \texttt{NmzConeProperty( cone, "NumberLatticePoints" );} see \texttt{NmzConeProperty} (\ref{NmzConeProperty}). }

 

\subsection{\textcolor{Chapter }{NmzOriginalMonoidGenerators}}
\logpage{[ 2, 3, 88 ]}\nobreak
\hyperdef{L}{X87F31DFA823B6D55}{}
{\noindent\textcolor{FuncColor}{$\triangleright$\enspace\texttt{NmzOriginalMonoidGenerators({\mdseries\slshape cone})\index{NmzOriginalMonoidGenerators@\texttt{NmzOriginalMonoidGenerators}}
\label{NmzOriginalMonoidGenerators}
}\hfill{\scriptsize (function)}}\\
\textbf{\indent Returns:\ }
a matrix whose rows are the original monoid generators 



 This is an alias for \texttt{NmzConeProperty( cone, "OriginalMonoidGenerators" );} see \texttt{NmzConeProperty} (\ref{NmzConeProperty}). }

 

\subsection{\textcolor{Chapter }{NmzPlacingTriangulation}}
\logpage{[ 2, 3, 89 ]}\nobreak
\hyperdef{L}{X7E9C4CC7878CF096}{}
{\noindent\textcolor{FuncColor}{$\triangleright$\enspace\texttt{NmzPlacingTriangulation({\mdseries\slshape cone})\index{NmzPlacingTriangulation@\texttt{NmzPlacingTriangulation}}
\label{NmzPlacingTriangulation}
}\hfill{\scriptsize (function)}}\\


 This is an alias for \texttt{NmzConeProperty( cone, "PlacingTriangulation" );} see \texttt{NmzConeProperty} (\ref{NmzConeProperty}). }

 

\subsection{\textcolor{Chapter }{NmzPrimalMode}}
\logpage{[ 2, 3, 90 ]}\nobreak
\hyperdef{L}{X7F201AE47FF7AEC6}{}
{\noindent\textcolor{FuncColor}{$\triangleright$\enspace\texttt{NmzPrimalMode({\mdseries\slshape cone})\index{NmzPrimalMode@\texttt{NmzPrimalMode}}
\label{NmzPrimalMode}
}\hfill{\scriptsize (function)}}\\


 This is an alias for \texttt{NmzConeProperty( cone, "PrimalMode" );} see \texttt{NmzConeProperty} (\ref{NmzConeProperty}). }

 

\subsection{\textcolor{Chapter }{NmzProjectCone}}
\logpage{[ 2, 3, 91 ]}\nobreak
\hyperdef{L}{X868659657E1463E3}{}
{\noindent\textcolor{FuncColor}{$\triangleright$\enspace\texttt{NmzProjectCone({\mdseries\slshape cone})\index{NmzProjectCone@\texttt{NmzProjectCone}}
\label{NmzProjectCone}
}\hfill{\scriptsize (function)}}\\


 

 Supported in Normaliz {\textgreater}= 3.5.0. 

 This is an alias for \texttt{NmzConeProperty( cone, "ProjectCone" );} see \texttt{NmzConeProperty} (\ref{NmzConeProperty}). }

 

\subsection{\textcolor{Chapter }{NmzProjection}}
\logpage{[ 2, 3, 92 ]}\nobreak
\hyperdef{L}{X7B5B2286812228B1}{}
{\noindent\textcolor{FuncColor}{$\triangleright$\enspace\texttt{NmzProjection({\mdseries\slshape cone})\index{NmzProjection@\texttt{NmzProjection}}
\label{NmzProjection}
}\hfill{\scriptsize (function)}}\\


 This is an alias for \texttt{NmzConeProperty( cone, "Projection" );} see \texttt{NmzConeProperty} (\ref{NmzConeProperty}). }

 

\subsection{\textcolor{Chapter }{NmzProjectionFloat}}
\logpage{[ 2, 3, 93 ]}\nobreak
\hyperdef{L}{X7913B00279250F90}{}
{\noindent\textcolor{FuncColor}{$\triangleright$\enspace\texttt{NmzProjectionFloat({\mdseries\slshape cone})\index{NmzProjectionFloat@\texttt{NmzProjectionFloat}}
\label{NmzProjectionFloat}
}\hfill{\scriptsize (function)}}\\


 This is an alias for \texttt{NmzConeProperty( cone, "ProjectionFloat" );} see \texttt{NmzConeProperty} (\ref{NmzConeProperty}). }

 

\subsection{\textcolor{Chapter }{NmzPullingTriangulation}}
\logpage{[ 2, 3, 94 ]}\nobreak
\hyperdef{L}{X811EC96985FE99FA}{}
{\noindent\textcolor{FuncColor}{$\triangleright$\enspace\texttt{NmzPullingTriangulation({\mdseries\slshape cone})\index{NmzPullingTriangulation@\texttt{NmzPullingTriangulation}}
\label{NmzPullingTriangulation}
}\hfill{\scriptsize (function)}}\\


 This is an alias for \texttt{NmzConeProperty( cone, "PullingTriangulation" );} see \texttt{NmzConeProperty} (\ref{NmzConeProperty}). }

 

\subsection{\textcolor{Chapter }{NmzPullingTriangulationInternal}}
\logpage{[ 2, 3, 95 ]}\nobreak
\hyperdef{L}{X8372059D783BA911}{}
{\noindent\textcolor{FuncColor}{$\triangleright$\enspace\texttt{NmzPullingTriangulationInternal({\mdseries\slshape cone})\index{NmzPullingTriangulationInternal@\texttt{NmzPullingTriangulationInternal}}
\label{NmzPullingTriangulationInternal}
}\hfill{\scriptsize (function)}}\\


 This is an alias for \texttt{NmzConeProperty( cone, "PullingTriangulationInternal" );} see \texttt{NmzConeProperty} (\ref{NmzConeProperty}). }

 

\subsection{\textcolor{Chapter }{NmzRank}}
\logpage{[ 2, 3, 96 ]}\nobreak
\hyperdef{L}{X7DE5D105856F5001}{}
{\noindent\textcolor{FuncColor}{$\triangleright$\enspace\texttt{NmzRank({\mdseries\slshape cone})\index{NmzRank@\texttt{NmzRank}}
\label{NmzRank}
}\hfill{\scriptsize (function)}}\\
\textbf{\indent Returns:\ }
the rank of the cone 



 

 This value is the rank of the lattice generated by the lattice points of the
cone. 

 This is an alias for \texttt{NmzConeProperty( cone, "Rank" );} see \texttt{NmzConeProperty} (\ref{NmzConeProperty}). }

 

\subsection{\textcolor{Chapter }{NmzRationalAutomorphisms}}
\logpage{[ 2, 3, 97 ]}\nobreak
\hyperdef{L}{X7982136F865B148A}{}
{\noindent\textcolor{FuncColor}{$\triangleright$\enspace\texttt{NmzRationalAutomorphisms({\mdseries\slshape cone})\index{NmzRationalAutomorphisms@\texttt{NmzRationalAutomorphisms}}
\label{NmzRationalAutomorphisms}
}\hfill{\scriptsize (function)}}\\


 This is an alias for \texttt{NmzConeProperty( cone, "RationalAutomorphisms" );} see \texttt{NmzConeProperty} (\ref{NmzConeProperty}). }

 

\subsection{\textcolor{Chapter }{NmzRecessionRank}}
\logpage{[ 2, 3, 98 ]}\nobreak
\hyperdef{L}{X7EE7F2F6841E3390}{}
{\noindent\textcolor{FuncColor}{$\triangleright$\enspace\texttt{NmzRecessionRank({\mdseries\slshape cone})\index{NmzRecessionRank@\texttt{NmzRecessionRank}}
\label{NmzRecessionRank}
}\hfill{\scriptsize (function)}}\\
\textbf{\indent Returns:\ }
the rank of the recession cone 



 

 Only for inhomogeneous computations. 

 This is an alias for \texttt{NmzConeProperty( cone, "RecessionRank" );} see \texttt{NmzConeProperty} (\ref{NmzConeProperty}). }

 

\subsection{\textcolor{Chapter }{NmzReesPrimaryMultiplicity}}
\logpage{[ 2, 3, 99 ]}\nobreak
\hyperdef{L}{X860FDA0D81BF641C}{}
{\noindent\textcolor{FuncColor}{$\triangleright$\enspace\texttt{NmzReesPrimaryMultiplicity({\mdseries\slshape cone})\index{NmzReesPrimaryMultiplicity@\texttt{NmzReesPrimaryMultiplicity}}
\label{NmzReesPrimaryMultiplicity}
}\hfill{\scriptsize (function)}}\\


 

 the multiplicity of a monomial ideal, provided it is primary to the maximal
ideal generated by the indeterminates. Used only with the input type \texttt{rees{\textunderscore}algebra}. 

 This is an alias for \texttt{NmzConeProperty( cone, "ReesPrimaryMultiplicity" );} see \texttt{NmzConeProperty} (\ref{NmzConeProperty}). }

 

\subsection{\textcolor{Chapter }{NmzRenfVolume}}
\logpage{[ 2, 3, 100 ]}\nobreak
\hyperdef{L}{X8172BBE57E970C22}{}
{\noindent\textcolor{FuncColor}{$\triangleright$\enspace\texttt{NmzRenfVolume({\mdseries\slshape cone})\index{NmzRenfVolume@\texttt{NmzRenfVolume}}
\label{NmzRenfVolume}
}\hfill{\scriptsize (function)}}\\


 This is an alias for \texttt{NmzConeProperty( cone, "RenfVolume" );} see \texttt{NmzConeProperty} (\ref{NmzConeProperty}). }

 

\subsection{\textcolor{Chapter }{NmzSignedDec}}
\logpage{[ 2, 3, 101 ]}\nobreak
\hyperdef{L}{X7F637D5F86E4B500}{}
{\noindent\textcolor{FuncColor}{$\triangleright$\enspace\texttt{NmzSignedDec({\mdseries\slshape cone})\index{NmzSignedDec@\texttt{NmzSignedDec}}
\label{NmzSignedDec}
}\hfill{\scriptsize (function)}}\\


 This is an alias for \texttt{NmzConeProperty( cone, "SignedDec" );} see \texttt{NmzConeProperty} (\ref{NmzConeProperty}). }

 

\subsection{\textcolor{Chapter }{NmzStanleyDec}}
\logpage{[ 2, 3, 102 ]}\nobreak
\hyperdef{L}{X81AF9A717BEB6B83}{}
{\noindent\textcolor{FuncColor}{$\triangleright$\enspace\texttt{NmzStanleyDec({\mdseries\slshape cone})\index{NmzStanleyDec@\texttt{NmzStanleyDec}}
\label{NmzStanleyDec}
}\hfill{\scriptsize (function)}}\\


 This is an alias for \texttt{NmzConeProperty( cone, "StanleyDec" );} see \texttt{NmzConeProperty} (\ref{NmzConeProperty}). }

 

\subsection{\textcolor{Chapter }{NmzStatic}}
\logpage{[ 2, 3, 103 ]}\nobreak
\hyperdef{L}{X86C5550E7ED226A2}{}
{\noindent\textcolor{FuncColor}{$\triangleright$\enspace\texttt{NmzStatic({\mdseries\slshape cone})\index{NmzStatic@\texttt{NmzStatic}}
\label{NmzStatic}
}\hfill{\scriptsize (function)}}\\


 This is an alias for \texttt{NmzConeProperty( cone, "Static" );} see \texttt{NmzConeProperty} (\ref{NmzConeProperty}). }

 

\subsection{\textcolor{Chapter }{NmzStrictIsoTypeCheck}}
\logpage{[ 2, 3, 104 ]}\nobreak
\hyperdef{L}{X7F7C052E805BCE42}{}
{\noindent\textcolor{FuncColor}{$\triangleright$\enspace\texttt{NmzStrictIsoTypeCheck({\mdseries\slshape cone})\index{NmzStrictIsoTypeCheck@\texttt{NmzStrictIsoTypeCheck}}
\label{NmzStrictIsoTypeCheck}
}\hfill{\scriptsize (function)}}\\


 This is an alias for \texttt{NmzConeProperty( cone, "StrictIsoTypeCheck" );} see \texttt{NmzConeProperty} (\ref{NmzConeProperty}). }

 

\subsection{\textcolor{Chapter }{NmzSublattice}}
\logpage{[ 2, 3, 105 ]}\nobreak
\hyperdef{L}{X7B0074147C3D10D6}{}
{\noindent\textcolor{FuncColor}{$\triangleright$\enspace\texttt{NmzSublattice({\mdseries\slshape cone})\index{NmzSublattice@\texttt{NmzSublattice}}
\label{NmzSublattice}
}\hfill{\scriptsize (function)}}\\


 This is an alias for \texttt{NmzConeProperty( cone, "Sublattice" );} see \texttt{NmzConeProperty} (\ref{NmzConeProperty}). }

 

\subsection{\textcolor{Chapter }{NmzSuppHypsFloat}}
\logpage{[ 2, 3, 106 ]}\nobreak
\hyperdef{L}{X870DD7FB85D80425}{}
{\noindent\textcolor{FuncColor}{$\triangleright$\enspace\texttt{NmzSuppHypsFloat({\mdseries\slshape cone})\index{NmzSuppHypsFloat@\texttt{NmzSuppHypsFloat}}
\label{NmzSuppHypsFloat}
}\hfill{\scriptsize (function)}}\\


 

 Supported in Normaliz {\textgreater}= 3.5.2. 

 This is an alias for \texttt{NmzConeProperty( cone, "SuppHypsFloat" );} see \texttt{NmzConeProperty} (\ref{NmzConeProperty}). }

 

\subsection{\textcolor{Chapter }{NmzSupportHyperplanes}}
\logpage{[ 2, 3, 107 ]}\nobreak
\hyperdef{L}{X87E51DB27D9007B1}{}
{\noindent\textcolor{FuncColor}{$\triangleright$\enspace\texttt{NmzSupportHyperplanes({\mdseries\slshape cone})\index{NmzSupportHyperplanes@\texttt{NmzSupportHyperplanes}}
\label{NmzSupportHyperplanes}
}\hfill{\scriptsize (function)}}\\
\textbf{\indent Returns:\ }
a matrix whose rows represent the support hyperplanes 



 

 The equations cut out the linear space generated by the cone. The equations,
congruences and support hyperplanes together describe the lattice points of
the cone. 

 This is an alias for \texttt{NmzConeProperty( cone, "SupportHyperplanes" );} see \texttt{NmzConeProperty} (\ref{NmzConeProperty}). }

 

\subsection{\textcolor{Chapter }{NmzSymmetrize}}
\logpage{[ 2, 3, 108 ]}\nobreak
\hyperdef{L}{X80C9E1F9849E5F0D}{}
{\noindent\textcolor{FuncColor}{$\triangleright$\enspace\texttt{NmzSymmetrize({\mdseries\slshape cone})\index{NmzSymmetrize@\texttt{NmzSymmetrize}}
\label{NmzSymmetrize}
}\hfill{\scriptsize (function)}}\\


 This is an alias for \texttt{NmzConeProperty( cone, "Symmetrize" );} see \texttt{NmzConeProperty} (\ref{NmzConeProperty}). }

 

\subsection{\textcolor{Chapter }{NmzTestArithOverflowDescent}}
\logpage{[ 2, 3, 109 ]}\nobreak
\hyperdef{L}{X85F2376580CED5C2}{}
{\noindent\textcolor{FuncColor}{$\triangleright$\enspace\texttt{NmzTestArithOverflowDescent({\mdseries\slshape cone})\index{NmzTestArithOverflowDescent@\texttt{NmzTestArithOverflowDescent}}
\label{NmzTestArithOverflowDescent}
}\hfill{\scriptsize (function)}}\\


 This is an alias for \texttt{NmzConeProperty( cone, "TestArithOverflowDescent" );} see \texttt{NmzConeProperty} (\ref{NmzConeProperty}). }

 

\subsection{\textcolor{Chapter }{NmzTestArithOverflowDualMode}}
\logpage{[ 2, 3, 110 ]}\nobreak
\hyperdef{L}{X804CF2FC7CFC9367}{}
{\noindent\textcolor{FuncColor}{$\triangleright$\enspace\texttt{NmzTestArithOverflowDualMode({\mdseries\slshape cone})\index{NmzTestArithOverflowDualMode@\texttt{NmzTestArithOverflowDualMode}}
\label{NmzTestArithOverflowDualMode}
}\hfill{\scriptsize (function)}}\\


 This is an alias for \texttt{NmzConeProperty( cone, "TestArithOverflowDualMode" );} see \texttt{NmzConeProperty} (\ref{NmzConeProperty}). }

 

\subsection{\textcolor{Chapter }{NmzTestArithOverflowFullCone}}
\logpage{[ 2, 3, 111 ]}\nobreak
\hyperdef{L}{X7E044A4D7E8D1207}{}
{\noindent\textcolor{FuncColor}{$\triangleright$\enspace\texttt{NmzTestArithOverflowFullCone({\mdseries\slshape cone})\index{NmzTestArithOverflowFullCone@\texttt{NmzTestArithOverflowFullCone}}
\label{NmzTestArithOverflowFullCone}
}\hfill{\scriptsize (function)}}\\


 This is an alias for \texttt{NmzConeProperty( cone, "TestArithOverflowFullCone" );} see \texttt{NmzConeProperty} (\ref{NmzConeProperty}). }

 

\subsection{\textcolor{Chapter }{NmzTestArithOverflowProjAndLift}}
\logpage{[ 2, 3, 112 ]}\nobreak
\hyperdef{L}{X82EC0B237F742AC9}{}
{\noindent\textcolor{FuncColor}{$\triangleright$\enspace\texttt{NmzTestArithOverflowProjAndLift({\mdseries\slshape cone})\index{NmzTestArithOverflowProjAndLift@\texttt{NmzTestArithOverflowProjAndLift}}
\label{NmzTestArithOverflowProjAndLift}
}\hfill{\scriptsize (function)}}\\


 This is an alias for \texttt{NmzConeProperty( cone, "TestArithOverflowProjAndLift" );} see \texttt{NmzConeProperty} (\ref{NmzConeProperty}). }

 

\subsection{\textcolor{Chapter }{NmzTestLargePyramids}}
\logpage{[ 2, 3, 113 ]}\nobreak
\hyperdef{L}{X7FEAF6A278552948}{}
{\noindent\textcolor{FuncColor}{$\triangleright$\enspace\texttt{NmzTestLargePyramids({\mdseries\slshape cone})\index{NmzTestLargePyramids@\texttt{NmzTestLargePyramids}}
\label{NmzTestLargePyramids}
}\hfill{\scriptsize (function)}}\\


 This is an alias for \texttt{NmzConeProperty( cone, "TestLargePyramids" );} see \texttt{NmzConeProperty} (\ref{NmzConeProperty}). }

 

\subsection{\textcolor{Chapter }{NmzTestLibNormaliz}}
\logpage{[ 2, 3, 114 ]}\nobreak
\hyperdef{L}{X7B6557CF84E93CA0}{}
{\noindent\textcolor{FuncColor}{$\triangleright$\enspace\texttt{NmzTestLibNormaliz({\mdseries\slshape cone})\index{NmzTestLibNormaliz@\texttt{NmzTestLibNormaliz}}
\label{NmzTestLibNormaliz}
}\hfill{\scriptsize (function)}}\\


 This is an alias for \texttt{NmzConeProperty( cone, "TestLibNormaliz" );} see \texttt{NmzConeProperty} (\ref{NmzConeProperty}). }

 

\subsection{\textcolor{Chapter }{NmzTestLinearAlgebraGMP}}
\logpage{[ 2, 3, 115 ]}\nobreak
\hyperdef{L}{X84412ABE84D42065}{}
{\noindent\textcolor{FuncColor}{$\triangleright$\enspace\texttt{NmzTestLinearAlgebraGMP({\mdseries\slshape cone})\index{NmzTestLinearAlgebraGMP@\texttt{NmzTestLinearAlgebraGMP}}
\label{NmzTestLinearAlgebraGMP}
}\hfill{\scriptsize (function)}}\\


 This is an alias for \texttt{NmzConeProperty( cone, "TestLinearAlgebraGMP" );} see \texttt{NmzConeProperty} (\ref{NmzConeProperty}). }

 

\subsection{\textcolor{Chapter }{NmzTestSimplexParallel}}
\logpage{[ 2, 3, 116 ]}\nobreak
\hyperdef{L}{X7CD00A81833E1899}{}
{\noindent\textcolor{FuncColor}{$\triangleright$\enspace\texttt{NmzTestSimplexParallel({\mdseries\slshape cone})\index{NmzTestSimplexParallel@\texttt{NmzTestSimplexParallel}}
\label{NmzTestSimplexParallel}
}\hfill{\scriptsize (function)}}\\


 This is an alias for \texttt{NmzConeProperty( cone, "TestSimplexParallel" );} see \texttt{NmzConeProperty} (\ref{NmzConeProperty}). }

 

\subsection{\textcolor{Chapter }{NmzTestSmallPyramids}}
\logpage{[ 2, 3, 117 ]}\nobreak
\hyperdef{L}{X80B85BB27B4ECC32}{}
{\noindent\textcolor{FuncColor}{$\triangleright$\enspace\texttt{NmzTestSmallPyramids({\mdseries\slshape cone})\index{NmzTestSmallPyramids@\texttt{NmzTestSmallPyramids}}
\label{NmzTestSmallPyramids}
}\hfill{\scriptsize (function)}}\\


 This is an alias for \texttt{NmzConeProperty( cone, "TestSmallPyramids" );} see \texttt{NmzConeProperty} (\ref{NmzConeProperty}). }

 

\subsection{\textcolor{Chapter }{NmzTriangulation}}
\logpage{[ 2, 3, 118 ]}\nobreak
\hyperdef{L}{X79E934EA873F442B}{}
{\noindent\textcolor{FuncColor}{$\triangleright$\enspace\texttt{NmzTriangulation({\mdseries\slshape cone})\index{NmzTriangulation@\texttt{NmzTriangulation}}
\label{NmzTriangulation}
}\hfill{\scriptsize (function)}}\\
\textbf{\indent Returns:\ }
the triangulation 



 

 This returns a list of the maximal simplicial cones in a triangulation, i.e.,
a list of cones dividing the cone into simplicial cones. Each cone in the list
is represented by a pair. The first entry of such a pair is the key of the
simplex, i.e., a list of integers $a_1,\dots,a_n$ referring to the \texttt{NmzGenerators} (\ref{NmzGenerators}) (counting from 0) which are used in this simplicial cone. The second entry of
each pair in the list is the absolute value of the determinant of the
generator matrix of the simplicial cone. 

 This is an alias for \texttt{NmzConeProperty( cone, "Triangulation" );} see \texttt{NmzConeProperty} (\ref{NmzConeProperty}). }

 

\subsection{\textcolor{Chapter }{NmzTriangulationDetSum}}
\logpage{[ 2, 3, 119 ]}\nobreak
\hyperdef{L}{X85BE26AF83A5BD27}{}
{\noindent\textcolor{FuncColor}{$\triangleright$\enspace\texttt{NmzTriangulationDetSum({\mdseries\slshape cone})\index{NmzTriangulationDetSum@\texttt{NmzTriangulationDetSum}}
\label{NmzTriangulationDetSum}
}\hfill{\scriptsize (function)}}\\
\textbf{\indent Returns:\ }
sum of the absolute values of the determinants of the simplicial cones in the
used triangulation 



 This is an alias for \texttt{NmzConeProperty( cone, "TriangulationDetSum" );} see \texttt{NmzConeProperty} (\ref{NmzConeProperty}). }

 

\subsection{\textcolor{Chapter }{NmzTriangulationSize}}
\logpage{[ 2, 3, 120 ]}\nobreak
\hyperdef{L}{X869840DC816DF323}{}
{\noindent\textcolor{FuncColor}{$\triangleright$\enspace\texttt{NmzTriangulationSize({\mdseries\slshape cone})\index{NmzTriangulationSize@\texttt{NmzTriangulationSize}}
\label{NmzTriangulationSize}
}\hfill{\scriptsize (function)}}\\
\textbf{\indent Returns:\ }
the number of simplicial cones in the used triangulation 



 This is an alias for \texttt{NmzConeProperty( cone, "TriangulationSize" );} see \texttt{NmzConeProperty} (\ref{NmzConeProperty}). }

 

\subsection{\textcolor{Chapter }{NmzUnimodularTriangulation}}
\logpage{[ 2, 3, 121 ]}\nobreak
\hyperdef{L}{X7C5DC23D7BD1DEEB}{}
{\noindent\textcolor{FuncColor}{$\triangleright$\enspace\texttt{NmzUnimodularTriangulation({\mdseries\slshape cone})\index{NmzUnimodularTriangulation@\texttt{NmzUnimodularTriangulation}}
\label{NmzUnimodularTriangulation}
}\hfill{\scriptsize (function)}}\\


 This is an alias for \texttt{NmzConeProperty( cone, "UnimodularTriangulation" );} see \texttt{NmzConeProperty} (\ref{NmzConeProperty}). }

 

\subsection{\textcolor{Chapter }{NmzUnitGroupIndex}}
\logpage{[ 2, 3, 122 ]}\nobreak
\hyperdef{L}{X7CDAD83D7FE1D02A}{}
{\noindent\textcolor{FuncColor}{$\triangleright$\enspace\texttt{NmzUnitGroupIndex({\mdseries\slshape cone})\index{NmzUnitGroupIndex@\texttt{NmzUnitGroupIndex}}
\label{NmzUnitGroupIndex}
}\hfill{\scriptsize (function)}}\\


 This is an alias for \texttt{NmzConeProperty( cone, "UnitGroupIndex" );} see \texttt{NmzConeProperty} (\ref{NmzConeProperty}). }

 

\subsection{\textcolor{Chapter }{NmzVerticesFloat}}
\logpage{[ 2, 3, 123 ]}\nobreak
\hyperdef{L}{X7FF8BB22787FA222}{}
{\noindent\textcolor{FuncColor}{$\triangleright$\enspace\texttt{NmzVerticesFloat({\mdseries\slshape cone})\index{NmzVerticesFloat@\texttt{NmzVerticesFloat}}
\label{NmzVerticesFloat}
}\hfill{\scriptsize (function)}}\\
\textbf{\indent Returns:\ }
a matrix whose rows are the vertices of the polyhedron \mbox{\texttt{\mdseries\slshape cone}} with float coordinates 



 

 The rows of this matrix represent the vertices of \mbox{\texttt{\mdseries\slshape cone}}, printed as floats for better readability. The result might be inexact, and
should therefore not be used for computations. 

 This is an alias for \texttt{NmzConeProperty( cone, "VerticesFloat" );} see \texttt{NmzConeProperty} (\ref{NmzConeProperty}). }

 

\subsection{\textcolor{Chapter }{NmzVerticesOfPolyhedron}}
\logpage{[ 2, 3, 124 ]}\nobreak
\hyperdef{L}{X84AB87777844250D}{}
{\noindent\textcolor{FuncColor}{$\triangleright$\enspace\texttt{NmzVerticesOfPolyhedron({\mdseries\slshape cone})\index{NmzVerticesOfPolyhedron@\texttt{NmzVerticesOfPolyhedron}}
\label{NmzVerticesOfPolyhedron}
}\hfill{\scriptsize (function)}}\\
\textbf{\indent Returns:\ }
a matrix whose rows are the vertices of the polyhedron 



 This is an alias for \texttt{NmzConeProperty( cone, "VerticesOfPolyhedron" );} see \texttt{NmzConeProperty} (\ref{NmzConeProperty}). }

 

\subsection{\textcolor{Chapter }{NmzVirtualMultiplicity}}
\logpage{[ 2, 3, 125 ]}\nobreak
\hyperdef{L}{X878AE417839ACF89}{}
{\noindent\textcolor{FuncColor}{$\triangleright$\enspace\texttt{NmzVirtualMultiplicity({\mdseries\slshape cone})\index{NmzVirtualMultiplicity@\texttt{NmzVirtualMultiplicity}}
\label{NmzVirtualMultiplicity}
}\hfill{\scriptsize (function)}}\\


 This is an alias for \texttt{NmzConeProperty( cone, "VirtualMultiplicity" );} see \texttt{NmzConeProperty} (\ref{NmzConeProperty}). }

 

\subsection{\textcolor{Chapter }{NmzVolume}}
\logpage{[ 2, 3, 126 ]}\nobreak
\hyperdef{L}{X8550FB5182075539}{}
{\noindent\textcolor{FuncColor}{$\triangleright$\enspace\texttt{NmzVolume({\mdseries\slshape cone})\index{NmzVolume@\texttt{NmzVolume}}
\label{NmzVolume}
}\hfill{\scriptsize (function)}}\\


 

 Supported in Normaliz {\textgreater}= 3.5.0. 

 This is an alias for \texttt{NmzConeProperty( cone, "Volume" );} see \texttt{NmzConeProperty} (\ref{NmzConeProperty}). }

 

\subsection{\textcolor{Chapter }{NmzWeightedEhrhartQuasiPolynomial}}
\logpage{[ 2, 3, 127 ]}\nobreak
\hyperdef{L}{X873956267E1EB3BD}{}
{\noindent\textcolor{FuncColor}{$\triangleright$\enspace\texttt{NmzWeightedEhrhartQuasiPolynomial({\mdseries\slshape cone})\index{NmzWeightedEhrhartQuasiPolynomial@\texttt{NmzWeightedEhrhartQuasiPolynomial}}
\label{NmzWeightedEhrhartQuasiPolynomial}
}\hfill{\scriptsize (function)}}\\


 This is an alias for \texttt{NmzConeProperty( cone, "WeightedEhrhartQuasiPolynomial" );} see \texttt{NmzConeProperty} (\ref{NmzConeProperty}). }

 

\subsection{\textcolor{Chapter }{NmzWeightedEhrhartSeries}}
\logpage{[ 2, 3, 128 ]}\nobreak
\hyperdef{L}{X7F7C11FD7C996E74}{}
{\noindent\textcolor{FuncColor}{$\triangleright$\enspace\texttt{NmzWeightedEhrhartSeries({\mdseries\slshape cone})\index{NmzWeightedEhrhartSeries@\texttt{NmzWeightedEhrhartSeries}}
\label{NmzWeightedEhrhartSeries}
}\hfill{\scriptsize (function)}}\\


 This is an alias for \texttt{NmzConeProperty( cone, "WeightedEhrhartSeries" );} see \texttt{NmzConeProperty} (\ref{NmzConeProperty}). }

 

\subsection{\textcolor{Chapter }{NmzWitnessNotIntegrallyClosed}}
\logpage{[ 2, 3, 129 ]}\nobreak
\hyperdef{L}{X7944885884E7368D}{}
{\noindent\textcolor{FuncColor}{$\triangleright$\enspace\texttt{NmzWitnessNotIntegrallyClosed({\mdseries\slshape cone})\index{NmzWitnessNotIntegrallyClosed@\texttt{NmzWitnessNotIntegrallyClosed}}
\label{NmzWitnessNotIntegrallyClosed}
}\hfill{\scriptsize (function)}}\\


 This is an alias for \texttt{NmzConeProperty( cone, "WitnessNotIntegrallyClosed" );} see \texttt{NmzConeProperty} (\ref{NmzConeProperty}). }

 

\subsection{\textcolor{Chapter }{NmzBasisChange}}
\logpage{[ 2, 3, 130 ]}\nobreak
\hyperdef{L}{X84D0AADB7BAF0328}{}
{\noindent\textcolor{FuncColor}{$\triangleright$\enspace\texttt{NmzBasisChange({\mdseries\slshape cone})\index{NmzBasisChange@\texttt{NmzBasisChange}}
\label{NmzBasisChange}
}\hfill{\scriptsize (function)}}\\
\textbf{\indent Returns:\ }
a record describing the basis change 



 The result record \texttt{r} has three components: \texttt{r.Embedding}, \texttt{r.Projection}, and \texttt{r.Annihilator}, where the embedding \texttt{A} and the projection \texttt{B} are matrices, and the annihilator \texttt{c} is an integer. They represent the mapping into the effective lattice $\mathbb{Z}^n \to \mathbb{Z}^r, u \mapsto (uB)/c$ and the inverse operation $\mathbb{Z}^r \to \mathbb{Z}^n, v \mapsto vA$. 

 This is part of the cone property ``Sublattice''. }

 }

 }

   
\chapter{\textcolor{Chapter }{Examples}}\label{Chapter_Examples}
\logpage{[ 3, 0, 0 ]}
\hyperdef{L}{X7A489A5D79DA9E5C}{}
{
  

 
\section{\textcolor{Chapter }{Generators}}\label{Chapter_Examples_Section_Generators}
\logpage{[ 3, 1, 0 ]}
\hyperdef{L}{X7BD5B55C802805B4}{}
{
  

 
\begin{Verbatim}[commandchars=!@|,fontsize=\small,frame=single,label=Example]
  !gapprompt@gap>| !gapinput@C := NmzCone(["integral_closure",[[2,1],[1,3]]]);|
  <a Normaliz cone>
  !gapprompt@gap>| !gapinput@NmzHasConeProperty(C,"HilbertBasis");|
  false
  !gapprompt@gap>| !gapinput@NmzHasConeProperty(C,"SupportHyperplanes");|
  false
  !gapprompt@gap>| !gapinput@NmzConeProperty(C,"HilbertBasis");|
  [ [ 1, 1 ], [ 1, 2 ], [ 1, 3 ], [ 2, 1 ] ]
  !gapprompt@gap>| !gapinput@NmzHasConeProperty(C,"SupportHyperplanes");|
  true
  !gapprompt@gap>| !gapinput@NmzConeProperty(C,"SupportHyperplanes");|
  [ [ -1, 2 ], [ 3, -1 ] ]
\end{Verbatim}
 

 }

 
\section{\textcolor{Chapter }{System of equations}}\label{Chapter_Examples_Section_System_of_equations}
\logpage{[ 3, 2, 0 ]}
\hyperdef{L}{X872A3FB37D6D7DFE}{}
{
  
\begin{Verbatim}[commandchars=!@|,fontsize=\small,frame=single,label=Example]
  !gapprompt@gap>| !gapinput@D := NmzCone(["equations",[[1,2,-3]], "grading",[[0,-1,3]]]);|
  <a Normaliz cone>
  !gapprompt@gap>| !gapinput@NmzCompute(D,["DualMode","HilbertSeries"]);|
  true
  !gapprompt@gap>| !gapinput@NmzHilbertBasis(D);|
  [ [ 1, 1, 1 ], [ 0, 3, 2 ], [ 3, 0, 1 ] ]
  !gapprompt@gap>| !gapinput@NmzHilbertSeries(D);|
  [ t^2-t+1, [ [ 1, 1 ], [ 3, 1 ] ] ]
  !gapprompt@gap>| !gapinput@NmzHasConeProperty(D,"SupportHyperplanes");|
  true
  !gapprompt@gap>| !gapinput@NmzSupportHyperplanes(D);|
  [ [ 0, 1, 0 ], [ 1, 0, 0 ] ]
  !gapprompt@gap>| !gapinput@NmzEquations(D);|
  [ [ 1, 2, -3 ] ]
\end{Verbatim}
 

 }

 
\section{\textcolor{Chapter }{System of inhomogeneous equations}}\label{Chapter_Examples_Section_System_of_inhomogeneous_equations}
\logpage{[ 3, 3, 0 ]}
\hyperdef{L}{X8570F3C87B5B7DD7}{}
{
  
\begin{Verbatim}[commandchars=!@|,fontsize=\small,frame=single,label=Example]
  !gapprompt@gap>| !gapinput@P := NmzCone(["inhom_equations",[[1,2,-3,1]], "grading", [[1,1,1]]]);|
  <a Normaliz cone>
  !gapprompt@gap>| !gapinput@NmzIsInhomogeneous(C);|
  false
  !gapprompt@gap>| !gapinput@NmzIsInhomogeneous(P);|
  true
  !gapprompt@gap>| !gapinput@NmzHilbertBasis(P);|
  [ [ 1, 1, 1, 0 ], [ 3, 0, 1, 0 ], [ 0, 3, 2, 0 ] ]
  !gapprompt@gap>| !gapinput@NmzModuleGenerators(P);|
  [ [ 0, 1, 1, 1 ], [ 2, 0, 1, 1 ] ]
\end{Verbatim}
 

 }

 
\section{\textcolor{Chapter }{Combined input}}\label{Chapter_Examples_Section_Combined_input}
\logpage{[ 3, 4, 0 ]}
\hyperdef{L}{X7B191A1778012C1D}{}
{
  Normaliz also allows the combination of different kinds of input, e.g.
multiple constraint types, or additional data like a grading. 

 Suppose that you have a 3 by 3 ``square'' of nonnegative integers such that the 3 numbers in all rows, all columns, and
both diagonals sum to the same constant $M$. Sometimes such squares are called magic and $M$ is the magic constant. This leads to a linear system of equations. The magic
constant is a natural choice for the grading. Additionally we force here the 4
corner to have even value by adding congruences. 

 
\begin{Verbatim}[commandchars=!@|,fontsize=\small,frame=single,label=Example]
  !gapprompt@gap>| !gapinput@Magic3x3even := NmzCone(["equations",|
  !gapprompt@>| !gapinput@[ [1, 1, 1,  -1, -1, -1,   0,  0,  0],|
  !gapprompt@>| !gapinput@  [1, 1, 1,   0,  0,  0,  -1, -1, -1],|
  !gapprompt@>| !gapinput@  [0, 1, 1,  -1,  0,  0,  -1,  0,  0],|
  !gapprompt@>| !gapinput@  [1, 0, 1,   0, -1,  0,   0, -1,  0],|
  !gapprompt@>| !gapinput@  [1, 1, 0,   0,  0, -1,   0,  0, -1],|
  !gapprompt@>| !gapinput@  [0, 1, 1,   0, -1,  0,   0,  0, -1],|
  !gapprompt@>| !gapinput@  [1, 1, 0,   0, -1,  0,  -1,  0,  0] ],|
  !gapprompt@>| !gapinput@"congruences",|
  !gapprompt@>| !gapinput@[ [1, 0, 0,   0, 0, 0,   0, 0, 0,  2],|
  !gapprompt@>| !gapinput@  [0, 0, 1,   0, 0, 0,   0, 0, 0,  2],|
  !gapprompt@>| !gapinput@  [0, 0, 0,   0, 0, 0,   1, 0, 0,  2],|
  !gapprompt@>| !gapinput@  [0, 0, 0,   0, 0, 0,   0, 0, 1,  2] ],|
  !gapprompt@>| !gapinput@"grading",|
  !gapprompt@>| !gapinput@[ [1, 1, 1,   0, 0, 0,   0, 0, 0] ] ] );|
  <a Normaliz cone>
  !gapprompt@gap>| !gapinput@NmzHilbertBasis(Magic3x3even);|
  [ [ 0, 4, 2, 4, 2, 0, 2, 0, 4 ], [ 2, 0, 4, 4, 2, 0, 0, 4, 2 ], 
    [ 2, 2, 2, 2, 2, 2, 2, 2, 2 ], [ 2, 4, 0, 0, 2, 4, 4, 0, 2 ], 
    [ 4, 0, 2, 0, 2, 4, 2, 4, 0 ], [ 2, 3, 4, 5, 3, 1, 2, 3, 4 ], 
    [ 2, 5, 2, 3, 3, 3, 4, 1, 4 ], [ 4, 1, 4, 3, 3, 3, 2, 5, 2 ], 
    [ 4, 3, 2, 1, 3, 5, 4, 3, 2 ] ]
  !gapprompt@gap>| !gapinput@NmzHilbertSeries(Magic3x3even);|
  [ t^3+3*t^2-t+1, [ [ 1, 1 ], [ 2, 2 ] ] ]
  !gapprompt@gap>| !gapinput@NmzHilbertQuasiPolynomial(Magic3x3even);|
  [ 1/2*t^2+t+1, 1/2*t^2-1/2 ]
\end{Verbatim}
 

 }

 
\section{\textcolor{Chapter }{Using the dual mode}}\label{Chapter_Examples_Section_Using_the_dual_mode}
\logpage{[ 3, 5, 0 ]}
\hyperdef{L}{X8406B01578B99D14}{}
{
  For solving systems of equations and inequalities it is often faster to use
the dual Normaliz algorithm. We demonstrate how to use it with an
inhomogeneous system of equations and inequalities. 

 The input consists of a system of 8 inhomogeneous equations in
R\texttt{\symbol{94}}3. The first row of the matrix M encodes the inequality $8x + 8y + 8z + 7 \geq 0$. Additionally we say that $x, y, z$ should be non\texttt{\symbol{45}}negative by giving the sign vector and use
the total degree. 
\begin{Verbatim}[commandchars=!@|,fontsize=\small,frame=single,label=Example]
  !gapprompt@gap>| !gapinput@M := [|
  !gapprompt@>| !gapinput@ [ 8,  8,  8,  7 ],|
  !gapprompt@>| !gapinput@ [ 0,  4,  0,  1 ],|
  !gapprompt@>| !gapinput@ [ 0,  1,  0,  7 ],|
  !gapprompt@>| !gapinput@ [ 0, -2,  0,  7 ],|
  !gapprompt@>| !gapinput@ [ 0, -2,  0,  1 ],|
  !gapprompt@>| !gapinput@ [ 8, 48,  8, 17 ],|
  !gapprompt@>| !gapinput@ [ 1,  6,  1, 34 ],|
  !gapprompt@>| !gapinput@ [ 2,-12, -2, 37 ],|
  !gapprompt@>| !gapinput@ [ 4,-24, -4, 14 ]|
  !gapprompt@>| !gapinput@];;|
  !gapprompt@gap>| !gapinput@D := NmzCone(["inhom_inequalities", M,|
  !gapprompt@>| !gapinput@              "signs", [[1,1,1]],|
  !gapprompt@>| !gapinput@              "grading", [[1,1,1]]]);|
  <a Normaliz cone>
  !gapprompt@gap>| !gapinput@NmzCompute(D,["DualMode","HilbertBasis","ModuleGenerators"]);|
  true
  !gapprompt@gap>| !gapinput@NmzHilbertBasis(D);|
  [ [ 1, 0, 0, 0 ], [ 1, 0, 1, 0 ] ]
  !gapprompt@gap>| !gapinput@NmzModuleGenerators(D);|
  [ [ 0, 0, 0, 1 ], [ 0, 0, 1, 1 ], [ 0, 0, 2, 1 ], [ 0, 0, 3, 1 ] ]
\end{Verbatim}
 

 As result we get the Hilbert basis of the cone of the solutions to the
homogeneous system and the module generators which are the base solutions to
the inhomogeneous system. 

 }

 }

   
\chapter{\textcolor{Chapter }{Installing NormalizInterface}}\label{Chapter_Installing_NormalizInterface}
\logpage{[ 4, 0, 0 ]}
\hyperdef{L}{X7DA4E7697F7D5F0C}{}
{
  

 
\section{\textcolor{Chapter }{Compiling}}\label{Chapter_Installing_NormalizInterface_Section_Compiling}
\logpage{[ 4, 1, 0 ]}
\hyperdef{L}{X7CD1A8937DFB78BF}{}
{
  

 NormalizInterface supports GAP 4.9 or later, and Normaliz 3.5.4 or later. 

 

 For technical reasons, installing and using NormalizInterface requires that
your version of GAP is compiled in a special way. Specifically, GAP must be
compiled against the exact same version of the GMP library as Normaliz. By
default, GAP compiles its own version of GMP; however, we cannot use that, as
it lacks C++ support, which is required by Normaliz. 

 

 Thus as the very first step, please install a version of GMP in your system.
On most Linux and BSD distributions, there should be a GMP package available
with your system's package manager. On Mac OS X, you can install GMP via Fink,
MacPorts or Homebrew. 

 

 Next, make sure your GAP installation is compiled against the system wide GMP
installation. To do so, switch to the GAP root directory, and enter the
following commands: 

 

 
\begin{Verbatim}[commandchars=!@|,fontsize=\small,frame=single,label=]
  make clean
  ./configure --with-gmp=PATH/TO/YOUR/GMP
  make
\end{Verbatim}
 

 Next you need to compile a recent version of Normaliz. This requires the
presence of several further system software packages, which you install via
your system's package manager. At least the following are required: 

 

 
\begin{itemize}
\item  curl OR wget for downloading the source code 
\end{itemize}
 

 Once you have installed these, you can build Normaliz by using the \texttt{prerequisites.sh} script we provide. It takes a single, optional parameter: the location of the
GAP root directory. 

 

 
\begin{Verbatim}[commandchars=!@|,fontsize=\small,frame=single,label=]
  ./prerequisites.sh GAPDIR
\end{Verbatim}
 

 Once it completed successfully, you can then build NormalizInterface like
this: 

 

 
\begin{Verbatim}[commandchars=!@|,fontsize=\small,frame=single,label=]
  ./configure --with-gaproot=GAPDIR
  make
\end{Verbatim}
 }

 }

 \def\bibname{References\logpage{[ "Bib", 0, 0 ]}
\hyperdef{L}{X7A6F98FD85F02BFE}{}
}

\bibliographystyle{alpha}
\bibliography{NormalizInterface-bib.xml}

\addcontentsline{toc}{chapter}{References}

\def\indexname{Index\logpage{[ "Ind", 0, 0 ]}
\hyperdef{L}{X83A0356F839C696F}{}
}

\cleardoublepage
\phantomsection
\addcontentsline{toc}{chapter}{Index}


\printindex

\immediate\write\pagenrlog{["Ind", 0, 0], \arabic{page},}
\newpage
\immediate\write\pagenrlog{["End"], \arabic{page}];}
\immediate\closeout\pagenrlog
\end{document}
